\documentclass[11pt,letterpaper]{report}
\usepackage[utf8]{inputenc}
\usepackage[margin=1in]{geometry}
\usepackage{amsmath}
\usepackage{graphicx}
\usepackage{hyperref}
\usepackage{array}
\usepackage{booktabs}
\usepackage{longtable}
\usepackage{fancyhdr}

\pagestyle{fancy}
\fancyhf{}
\fancyhead[L]{\leftmark}
\fancyfoot[C]{\thepage}
\fancyhead[R]{Exam Experts}
\renewcommand{\headrulewidth}{0.4pt}

\hypersetup{
    colorlinks=true,
    linkcolor=blue,
    urlcolor=blue,
    citecolor=blue
}

\title{\textbf{COMPREHENSIVE ELEMENTARY AND MIDDLE SCHOOL\\
COURSE SYLLABI\\
COMPLETE EDITION}\\
\vspace{0.5cm}
\large Institutional-Grade Curriculum Documentation\\
Complete Curriculum Guide for Kindergarten Through 8th Grade}
\author{}
\date{January 2026}

\begin{document}

\maketitle

\begin{center}
\textbf{Includes:}\\
\vspace{0.3cm}
\begin{itemize}
    \item Core Subjects: English/Language Arts, Mathematics, Science, Social Studies
    \item Special Subjects: Arts, Music, Physical Education, Health
    \item Complete Grade-by-Grade Curriculum Framework (K-8)
    \item Standards-Aligned Learning Objectives
    \item Differentiation Strategies for Diverse Learners
\end{itemize}

\vspace{0.5cm}

\textbf{Developed in accordance with:}\\
\vspace{0.3cm}
\begin{itemize}
    \item Common Core State Standards (CCSS)
    \item Next Generation Science Standards (NGSS)
    \item National Council for the Social Studies (NCSS) Standards
    \item National Standards for Arts Education
    \item SHAPE America Physical Education Standards
\end{itemize}
\end{center}

\clearpage

\tableofcontents

\clearpage

\chapter*{Preface}
\addcontentsline{toc}{chapter}{Preface}

This comprehensive curriculum documentation provides institutional-grade syllabi for all major elementary and middle school courses from Kindergarten through 8th Grade. Each syllabus has been developed to meet national and state standards while providing flexibility for local adaptation and differentiation to meet diverse student needs.

\section*{Purpose}

This document serves multiple audiences:

\begin{itemize}
    \item \textbf{Teachers:} Comprehensive curriculum framework for instruction across all grade levels
    \item \textbf{Administrators:} Documentation for accreditation, program evaluation, and curriculum alignment
    \item \textbf{Parents:} Understanding of grade-level expectations and content progression
    \item \textbf{Curriculum Coordinators:} Vertical alignment and scope and sequence planning
    \item \textbf{Student Teachers:} Professional preparation and understanding of developmental progressions
\end{itemize}

\section*{Organization}

This guide is organized by grade level (Kindergarten through 8th Grade), with each grade subdivided by subject area:

\begin{itemize}
    \item English/Language Arts
    \item Mathematics
    \item Science
    \item Social Studies
    \item Additional subjects (Arts, Music, PE, Health) covered at grade level
\end{itemize}

\section*{Standards Alignment}

All syllabi align with:

\begin{itemize}
    \item Common Core State Standards for English Language Arts and Mathematics
    \item Next Generation Science Standards (NGSS)
    \item National Council for the Social Studies (NCSS) C3 Framework
    \item National Core Arts Standards
    \item SHAPE America National Standards for Physical Education
    \item National Health Education Standards
\end{itemize}

\section*{Differentiation and Inclusion}

Throughout this guide, emphasis is placed on:

\begin{itemize}
    \item Differentiated instruction strategies
    \item Accommodations for English Language Learners (ELL)
    \item Modifications for students with special needs
    \item Enrichment opportunities for gifted learners
    \item Universal Design for Learning (UDL) principles
\end{itemize}

\section*{Assessment}

Assessment strategies include:

\begin{itemize}
    \item Formative assessments (ongoing checks for understanding)
    \item Summative assessments (unit tests, projects, performance tasks)
    \item Diagnostic assessments (pre-assessments, screening)
    \item Portfolio-based assessments
    \item Standards-based grading approaches
\end{itemize}

\clearpage

\part{Kindergarten}

\chapter{Kindergarten English/Language Arts}

\section{Course Overview}

\textbf{Grade Level:} Kindergarten\\
\textbf{Duration:} Full Year (36 weeks)\\
\textbf{Prerequisites:} None

\section{Course Description}

Kindergarten English/Language Arts focuses on developing foundational literacy skills including phonological awareness, letter recognition, early reading skills, listening comprehension, oral language development, and beginning writing. Students engage with stories, informational texts, and poetry while building vocabulary and communication skills essential for future academic success.

\section{Course Goals}

Upon completion of this course, students will be able to:

\begin{itemize}
    \item Demonstrate understanding of spoken words, syllables, and sounds (phonemes)
    \item Recognize and name all upper and lowercase letters
    \item Read common high-frequency words by sight
    \item Read emergent-reader texts with purpose and understanding
    \item Use a combination of drawing, dictating, and writing to compose texts
    \item Participate in collaborative conversations with peers and adults
    \item Ask and answer questions about key details in texts read aloud
    \item Describe familiar people, places, things, and events
\end{itemize}

\section{Units of Study}

\subsection{Unit 1: Print Concepts and Phonological Awareness (Ongoing)}

\textbf{Duration:} Integrated throughout the year

\textbf{Topics:}
\begin{itemize}
    \item Understanding that print carries meaning
    \item Following words from left to right, top to bottom, page by page
    \item Recognizing that spoken words are represented by written words
    \item Understanding that words are separated by spaces
    \item Recognizing and producing rhyming words
    \item Counting, pronouncing, blending, and segmenting syllables
    \item Blending and segmenting onsets and rimes
    \item Isolating and pronouncing initial, medial vowel, and final sounds (phonemes)
    \item Adding or substituting sounds to make new words
\end{itemize}

\subsection{Unit 2: Letter Recognition and Phonics (8-10 weeks)}

\textbf{Topics:}
\begin{itemize}
    \item Recognizing and naming all 26 uppercase letters
    \item Recognizing and naming all 26 lowercase letters
    \item Producing the primary sound for each consonant
    \item Associating long and short vowel sounds with common spellings
    \item Reading common high-frequency words by sight (e.g., the, of, to, you, she, my, is, are, do, does)
    \item Distinguishing between letters and non-letters
    \item Understanding alphabetical order
\end{itemize}

\subsection{Unit 3: Reading Foundation Skills (10-12 weeks)}

\textbf{Topics:}
\begin{itemize}
    \item Reading emergent-reader texts with purpose and understanding
    \item Demonstrating one-to-one correspondence (pointing to words while reading)
    \item Using picture cues and context to support reading
    \item Applying phonics skills to decode simple words
    \item Recognizing common sight words automatically
    \item Self-monitoring and self-correcting while reading
    \item Reading with expression and appropriate pacing
\end{itemize}

\subsection{Unit 4: Reading Literature (10-12 weeks)}

\textbf{Topics:}
\begin{itemize}
    \item Asking and answering questions about key details in stories
    \item Retelling familiar stories with key details
    \item Identifying characters, settings, and major events in stories
    \item Comparing and contrasting adventures and experiences of characters
    \item Actively engaging in group reading activities with purpose and understanding
    \item Recognizing common types of texts (storybooks, poems)
    \item Making connections between texts and personal experiences
\end{itemize}

\subsection{Unit 5: Reading Informational Text (8-10 weeks)}

\textbf{Topics:}
\begin{itemize}
    \item Asking and answering questions about key details in informational texts
    \item Identifying the main topic and retelling key details
    \item Describing the connection between two individuals, events, ideas, or pieces of information
    \item Identifying the front cover, back cover, and title page of a book
    \item Naming the author and illustrator and defining their roles
    \item Describing the relationship between illustrations and text
\end{itemize}

\subsection{Unit 6: Writing (Ongoing)}

\textbf{Duration:} Integrated throughout the year

\textbf{Topics:}
\begin{itemize}
    \item Using a combination of drawing, dictating, and writing to compose opinion pieces
    \item Using a combination of drawing, dictating, and writing to compose informative/explanatory texts
    \item Using a combination of drawing, dictating, and writing to narrate events
    \item Responding to questions and suggestions to strengthen writing
    \item Exploring digital tools to produce and publish writing with adult support
    \item Participating in shared research and writing projects
    \item Forming uppercase and lowercase letters correctly
    \item Spelling simple words phonetically
\end{itemize}

\subsection{Unit 7: Speaking and Listening (Ongoing)}

\textbf{Duration:} Integrated throughout the year

\textbf{Topics:}
\begin{itemize}
    \item Participating in collaborative conversations with diverse partners
    \item Following agreed-upon rules for discussions (listening, taking turns)
    \item Continuing conversations through multiple exchanges
    \item Asking and answering questions to seek help, get information, or clarify
    \item Describing familiar people, places, things, and events with detail
    \item Speaking audibly and expressing thoughts, feelings, and ideas clearly
\end{itemize}

\subsection{Unit 8: Language Skills (Ongoing)}

\textbf{Duration:} Integrated throughout the year

\textbf{Topics:}
\begin{itemize}
    \item Printing many upper and lowercase letters
    \item Using frequently occurring nouns and verbs
    \item Forming regular plural nouns orally
    \item Understanding and using question words (who, what, where, when, why, how)
    \item Using prepositions (to, from, in, out, on, off, for, of, by, with)
    \item Producing and expanding complete sentences
    \item Capitalizing the first word in a sentence and the pronoun I
    \item Recognizing and naming end punctuation
    \item Writing a letter or letters for most consonant and short-vowel sounds
    \item Determining or clarifying the meaning of unknown words
    \item Exploring word relationships and nuances in word meanings
\end{itemize}

\section{Required Materials}

\begin{itemize}
    \item Alphabet charts and letter cards
    \item Leveled readers (emergent and early levels)
    \item Picture books for read-alouds
    \item Writing journals or notebooks
    \item Manipulatives for phonemic awareness (counters, blocks)
    \item Classroom library with diverse texts
    \item Art supplies for drawing and writing
    \item Digital devices for technology integration (optional)
\end{itemize}

\section{Assessment and Grading}

Kindergarten assessment focuses on developmental progress rather than traditional grading:

\begin{center}
\begin{tabular}{|l|p{8cm}|}
\hline
\textbf{Assessment Type} & \textbf{Description} \\
\hline
Observational Notes & Daily observations of student engagement and skill application \\
\hline
Running Records & One-on-one reading assessments to monitor decoding and fluency \\
\hline
Portfolio Assessment & Collection of student work samples showing growth over time \\
\hline
Phonics Screeners & Assessment of letter knowledge and phonics skills \\
\hline
Conference Notes & Individual conferences to assess comprehension and set goals \\
\hline
\end{tabular}
\end{center}

\textbf{Progress Reports:}
\begin{itemize}
    \item Skills-based progress indicators (Beginning, Developing, Proficient)
    \item Narrative comments on student growth and areas for development
    \item Goal-setting for continued learning
\end{itemize}

\chapter{Kindergarten Mathematics}

\section{Course Overview}

\textbf{Grade Level:} Kindergarten\\
\textbf{Duration:} Full Year (36 weeks)\\
\textbf{Prerequisites:} None

\section{Course Description}

Kindergarten Mathematics develops foundational understanding of counting, number sense, basic operations, and geometric concepts. Students engage in hands-on activities, games, and problem-solving experiences that build mathematical thinking and reasoning skills. The curriculum emphasizes making sense of mathematics through exploration, discussion, and representation.

\section{Course Goals}

Upon completion of this course, students will be able to:

\begin{itemize}
    \item Count to 100 by ones and tens
    \item Count forward beginning from any given number (not just 1)
    \item Write numbers from 0 to 20
    \item Understand the relationship between numbers and quantities
    \item Compare numbers and understand greater than, less than, and equal to
    \item Represent addition and subtraction with objects, fingers, drawings, and equations
    \item Solve addition and subtraction word problems within 10
    \item Identify, describe, and create two-dimensional and three-dimensional shapes
    \item Compare measurable attributes of objects (longer/shorter, heavier/lighter)
\end{itemize}

\section{Units of Study}

\subsection{Unit 1: Math in Our World (3-4 weeks)}

\textbf{Topics:}
\begin{itemize}
    \item Exploring and using math tools (counters, pattern blocks, ten-frames)
    \item Recognizing math in everyday life
    \item Developing number sense through play and exploration
    \item Sorting and classifying objects by attributes
    \item Creating and extending simple patterns
    \item Introduction to mathematical language and vocabulary
\end{itemize}

\subsection{Unit 2: Numbers 1-10 (5-6 weeks)}

\textbf{Topics:}
\begin{itemize}
    \item Counting objects to 10
    \item Understanding one-to-one correspondence
    \item Writing numbers 0-10
    \item Recognizing numerals 0-10
    \item Comparing numbers within 10 (more, less, same)
    \item Understanding cardinality (the last number said tells how many)
    \item Counting to answer ``how many?'' questions
    \item Composing and decomposing numbers to 10
\end{itemize}

\subsection{Unit 3: Flat Shapes All Around Us (3-4 weeks)}

\textbf{Topics:}
\begin{itemize}
    \item Identifying and naming shapes (circle, triangle, square, rectangle, hexagon)
    \item Describing shapes using positional words (above, below, beside, in front of, behind, next to)
    \item Analyzing and comparing two-dimensional shapes
    \item Composing simple shapes to form larger shapes
    \item Recognizing shapes in the environment
    \item Drawing shapes
\end{itemize}

\subsection{Unit 4: Understanding Addition and Subtraction (4-5 weeks)}

\textbf{Topics:}
\begin{itemize}
    \item Representing addition as putting together
    \item Representing subtraction as taking apart or taking from
    \item Using objects, fingers, and drawings to model problems
    \item Solving addition and subtraction word problems within 10
    \item Understanding the concept of equals (=)
    \item Writing simple addition and subtraction equations
\end{itemize}

\subsection{Unit 5: Composing and Decomposing Numbers to 10 (4-5 weeks)}

\textbf{Topics:}
\begin{itemize}
    \item Finding all the ways to make a number (e.g., 5 = 4+1, 5 = 3+2)
    \item Using ten-frames to visualize number relationships
    \item Decomposing numbers in more than one way
    \item Understanding part-part-whole relationships
    \item Fluently adding and subtracting within 5
    \item Number bonds and fact families
\end{itemize}

\subsection{Unit 6: Numbers 0-20 (5-6 weeks)}

\textbf{Topics:}
\begin{itemize}
    \item Counting to 20 and beyond
    \item Writing numbers from 0 to 20
    \item Counting forward from any number within 20
    \item Understanding teen numbers as ten and some ones
    \item Comparing numbers within 20
    \item Ordering numbers from least to greatest
\end{itemize}

\subsection{Unit 7: Solid Shapes All Around Us (3-4 weeks)}

\textbf{Topics:}
\begin{itemize}
    \item Identifying and naming three-dimensional shapes (sphere, cube, cone, cylinder)
    \item Describing three-dimensional shapes (flat surfaces, curved surfaces, vertices)
    \item Comparing two-dimensional and three-dimensional shapes
    \item Building and creating with three-dimensional shapes
    \item Recognizing 3D shapes in the environment
\end{itemize}

\subsection{Unit 8: Measurement and Data (3-4 weeks)}

\textbf{Topics:}
\begin{itemize}
    \item Comparing objects by length (longer, shorter, same length)
    \item Comparing objects by height (taller, shorter, same height)
    \item Comparing objects by weight (heavier, lighter, same weight)
    \item Comparing objects by capacity (holds more, holds less, holds the same)
    \item Classifying objects into categories and counting
    \item Creating simple graphs and charts
\end{itemize}

\subsection{Unit 9: Putting It All Together (3-4 weeks)}

\textbf{Topics:}
\begin{itemize}
    \item Reviewing and reinforcing all concepts
    \item Applying multiple concepts to solve problems
    \item Engaging in mathematical games and centers
    \item Problem-solving with real-world contexts
    \item Celebrating mathematical growth and learning
\end{itemize}

\section{Required Materials}

\begin{itemize}
    \item Manipulatives: counting bears, unifix cubes, pattern blocks
    \item Ten-frames and number lines
    \item Number cards and dot cards
    \item Two-dimensional and three-dimensional shape sets
    \item Math games and puzzles
    \item Measurement tools (balance scale, measuring cups, linking cubes)
    \item Math journals or notebooks
    \item Classroom number line (0-20 or 0-100)
\end{itemize}

\section{Assessment and Grading}

\begin{center}
\begin{tabular}{|l|p{8cm}|}
\hline
\textbf{Assessment Type} & \textbf{Description} \\
\hline
Math Talks & Small group discussions to assess understanding \\
\hline
Performance Tasks & Hands-on activities demonstrating conceptual understanding \\
\hline
Observational Assessment & Daily observations during math centers and lessons \\
\hline
Individual Interviews & One-on-one assessments of counting, number sense, problem-solving \\
\hline
Portfolio Work & Collection of student work samples and problem-solving \\
\hline
\end{tabular}
\end{center}

\chapter{Kindergarten Science}

\section{Course Overview}

\textbf{Grade Level:} Kindergarten\\
\textbf{Duration:} Full Year (36 weeks)\\
\textbf{Prerequisites:} None

\section{Course Description}

Kindergarten Science introduces students to scientific inquiry through hands-on exploration and observation. Students investigate the natural world, asking questions and seeking answers through investigations. The curriculum integrates life science, physical science, earth and space science, and engineering design, fostering curiosity and developing foundational science practices.

\section{Course Goals}

Upon completion of this course, students will be able to:

\begin{itemize}
    \item Ask questions based on observations
    \item Plan and conduct simple investigations with guidance
    \item Make observations using the five senses
    \item Use tools to gather and organize data
    \item Construct explanations based on observations
    \item Communicate findings through drawings, writing, and speaking
    \item Understand that plants and animals have needs
    \item Observe weather patterns and seasonal changes
    \item Investigate properties of matter
    \item Explore pushes and pulls (forces)
\end{itemize}

\section{Units of Study}

\subsection{Unit 1: Science Skills and Inquiry (Ongoing)}

\textbf{Duration:} Integrated throughout the year

\textbf{Topics:}
\begin{itemize}
    \item Asking questions and making predictions
    \item Using the five senses to observe
    \item Using simple tools (hand lenses, measuring tools)
    \item Recording observations through drawings and writing
    \item Comparing and classifying objects
    \item Communicating findings with others
\end{itemize}

\subsection{Unit 2: Plants and Animals (6-7 weeks)}

\textbf{Topics:}
\begin{itemize}
    \item Identifying living and non-living things
    \item Understanding that plants need water, air, and light to grow
    \item Understanding that animals need food, water, and shelter
    \item Observing plant life cycles (seed to plant)
    \item Comparing different types of plants and animals
    \item Understanding basic animal needs and habitats
    \item Caring for living things responsibly
\end{itemize}

\textbf{Investigations:}
\begin{itemize}
    \item Growing plants from seeds (observing growth over time)
    \item Comparing plants with and without water or light
    \item Observing animals (classroom pets, insects, or videos)
    \item Creating animal habitats
\end{itemize}

\subsection{Unit 3: Weather and Seasons (5-6 weeks)}

\textbf{Topics:}
\begin{itemize}
    \item Observing daily weather patterns
    \item Recording weather data (sunny, cloudy, rainy, snowy, windy)
    \item Understanding the four seasons
    \item Recognizing seasonal changes in weather, plants, and animals
    \item Identifying appropriate clothing for different weather
    \item Understanding that the sun provides light and warmth
\end{itemize}

\textbf{Investigations:}
\begin{itemize}
    \item Daily weather observations and charting
    \item Creating weather graphs
    \item Seasonal nature walks and observations
    \item Exploring temperature with thermometers
\end{itemize}

\subsection{Unit 4: Properties of Matter (5-6 weeks)}

\textbf{Topics:}
\begin{itemize}
    \item Observing and describing properties of objects (color, size, shape, texture, weight)
    \item Comparing and classifying materials
    \item Understanding that objects are made from different materials
    \item Exploring solid and liquid states of matter
    \item Investigating how materials change (melting, mixing, dissolving)
\end{itemize}

\textbf{Investigations:}
\begin{itemize}
    \item Sorting objects by properties
    \item Testing materials (sink/float, magnetic/non-magnetic)
    \item Melting ice and observing changes
    \item Mixing colors and materials
\end{itemize}

\subsection{Unit 5: Pushes and Pulls (Forces) (4-5 weeks)}

\textbf{Topics:}
\begin{itemize}
    \item Understanding that pushes and pulls make things move
    \item Exploring how the strength of a push or pull affects motion
    \item Investigating how the direction of a push or pull affects motion
    \item Recognizing pushes and pulls in everyday life
    \item Understanding that objects at rest stay at rest until pushed or pulled
\end{itemize}

\textbf{Investigations:}
\begin{itemize}
    \item Testing different pushes and pulls with toy cars
    \item Exploring ramps and inclines
    \item Investigating magnets (pushing and pulling without touching)
    \item Creating simple machines or contraptions
\end{itemize}

\subsection{Unit 6: Earth's Resources (4-5 weeks)}

\textbf{Topics:}
\begin{itemize}
    \item Understanding that we get resources from the earth (water, soil, rocks)
    \item Exploring different types of soil and rocks
    \item Learning about water as an essential resource
    \item Understanding the importance of taking care of our environment
    \item Reducing, reusing, and recycling
\end{itemize}

\textbf{Investigations:}
\begin{itemize}
    \item Examining soil samples
    \item Observing and sorting rocks
    \item Exploring water properties
    \item Creating recycling projects
\end{itemize}

\subsection{Unit 7: The Human Body (4-5 weeks)}

\textbf{Topics:}
\begin{itemize}
    \item Identifying basic body parts and their functions
    \item Understanding the five senses
    \item Learning about healthy habits (exercise, nutrition, hygiene)
    \item Understanding that our bodies grow and change
    \item Recognizing similarities and differences among people
\end{itemize}

\textbf{Investigations:}
\begin{itemize}
    \item Exploring the five senses through activities
    \item Measuring height and comparing growth
    \item Testing sense of taste, smell, touch, hearing, sight
\end{itemize}

\section{Required Materials}

\begin{itemize}
    \item Science journals or observation notebooks
    \item Hand lenses and magnifying glasses
    \item Thermometers (large demonstration type)
    \item Balance scales
    \item Plant seeds and growing containers
    \item Weather observation charts
    \item Natural materials (rocks, soil, leaves, seeds)
    \item Magnets and magnetic materials
    \item Art supplies for science drawings
    \item Safety equipment (as needed for investigations)
\end{itemize}

\section{Assessment and Grading}

\begin{center}
\begin{tabular}{|l|p{8cm}|}
\hline
\textbf{Assessment Type} & \textbf{Description} \\
\hline
Science Journals & Drawings and writing to document observations \\
\hline
Performance Tasks & Hands-on investigations and demonstrations \\
\hline
Oral Explanations & Verbal communication of scientific thinking \\
\hline
Observational Notes & Teacher observations during inquiry activities \\
\hline
Science Projects & Simple projects demonstrating understanding \\
\hline
\end{tabular}
\end{center}

\chapter{Kindergarten Social Studies}

\section{Course Overview}

\textbf{Grade Level:} Kindergarten\\
\textbf{Duration:} Full Year (36 weeks)\\
\textbf{Prerequisites:} None

\section{Course Description}

Kindergarten Social Studies helps students understand themselves, their families, their school, and their community. Students develop an awareness of rules, roles, and responsibilities while exploring basic geography, economics, and cultural diversity. The curriculum emphasizes citizenship, community helpers, and understanding the world around us.

\section{Course Goals}

Upon completion of this course, students will be able to:

\begin{itemize}
    \item Understand their role as a member of a family, classroom, and community
    \item Recognize and follow rules and routines
    \item Identify community helpers and their roles
    \item Understand basic map skills and spatial relationships
    \item Recognize national symbols and traditions
    \item Understand basic economic concepts (wants and needs, goods and services)
    \item Appreciate cultural diversity and similarities
    \item Demonstrate good citizenship behaviors
\end{itemize}

\section{Units of Study}

\subsection{Unit 1: All About Me (3-4 weeks)}

\textbf{Topics:}
\begin{itemize}
    \item Learning about ourselves (name, age, interests)
    \item Understanding what makes us unique and special
    \item Identifying personal characteristics and preferences
    \item Comparing similarities and differences among classmates
    \item Understanding personal history (growing and changing)
\end{itemize}

\subsection{Unit 2: My Family (3-4 weeks)}

\textbf{Topics:}
\begin{itemize}
    \item Identifying family members and family structures
    \item Understanding family roles and responsibilities
    \item Recognizing different types of families
    \item Learning about family traditions and customs
    \item Understanding how families meet needs
\end{itemize}

\subsection{Unit 3: School and Classroom Community (4-5 weeks)}

\textbf{Topics:}
\begin{itemize}
    \item Understanding school rules and routines
    \item Identifying school staff and their roles
    \item Learning about classroom responsibilities
    \item Developing respect for others and their property
    \item Participating as a member of the classroom community
    \item Following directions and procedures
\end{itemize}

\subsection{Unit 4: Community and Community Helpers (5-6 weeks)}

\textbf{Topics:}
\begin{itemize}
    \item Defining what a community is
    \item Identifying community helpers (police, firefighters, doctors, teachers, etc.)
    \item Understanding the jobs and roles of community helpers
    \item Learning about places in our community
    \item Recognizing how community members help each other
    \item Understanding goods and services in the community
\end{itemize}

\subsection{Unit 5: Maps and Geography (4-5 weeks)}

\textbf{Topics:}
\begin{itemize}
    \item Understanding that maps show places
    \item Using simple maps (classroom, school, neighborhood)
    \item Understanding spatial relationships (near, far, left, right, above, below)
    \item Identifying land and water on maps and globes
    \item Creating simple maps
    \item Understanding addresses and locations
\end{itemize}

\subsection{Unit 6: Our Country and National Symbols (4-5 weeks)}

\textbf{Topics:}
\begin{itemize}
    \item Identifying the United States on a map
    \item Recognizing the American flag and what it represents
    \item Learning about national symbols (Statue of Liberty, bald eagle, etc.)
    \item Understanding patriotic songs and holidays
    \item Showing respect for the flag and country
\end{itemize}

\subsection{Unit 7: Cultures and Holidays (4-5 weeks)}

\textbf{Topics:}
\begin{itemize}
    \item Learning about different cultures and traditions
    \item Celebrating holidays and special days
    \item Understanding that families celebrate in different ways
    \item Appreciating diversity and similarities
    \item Sharing cultural traditions
\end{itemize}

\subsection{Unit 8: Wants, Needs, and Economics (3-4 weeks)}

\textbf{Topics:}
\begin{itemize}
    \item Distinguishing between wants and needs
    \item Understanding goods and services
    \item Learning about jobs and earning money
    \item Making choices about spending
    \item Understanding that resources are limited
\end{itemize}

\subsection{Unit 9: Citizenship and Character (Ongoing)}

\textbf{Duration:} Integrated throughout the year

\textbf{Topics:}
\begin{itemize}
    \item Demonstrating respect, responsibility, and honesty
    \item Showing kindness and fairness to others
    \item Following rules and routines
    \item Helping others and working cooperatively
    \item Making good choices
    \item Resolving conflicts peacefully
\end{itemize}

\section{Required Materials}

\begin{itemize}
    \item Picture books about families, communities, and cultures
    \item Simple maps (classroom, school, community)
    \item Globe and United States map
    \item American flag and national symbols posters
    \item Community helper props and costumes (for dramatic play)
    \item Cultural artifacts and photos
    \item Chart paper for creating class books
\end{itemize}

\section{Assessment and Grading}

\begin{center}
\begin{tabular}{|l|p{8cm}|}
\hline
\textbf{Assessment Type} & \textbf{Description} \\
\hline
Observations & Daily observations of citizenship behaviors \\
\hline
Discussions & Class conversations and sharing \\
\hline
Projects & Simple projects (family tree, community helper poster) \\
\hline
Drawings & Visual representations of concepts \\
\hline
Role-Playing & Dramatic play activities demonstrating understanding \\
\hline
\end{tabular}
\end{center}

\clearpage

\part{Grade 1}

\chapter{Grade 1 English/Language Arts}

\section{Course Overview}

\textbf{Grade Level:} Grade 1\\
\textbf{Duration:} Full Year (36 weeks)\\
\textbf{Prerequisites:} Kindergarten or demonstrated readiness

\section{Course Description}

Grade 1 English/Language Arts builds upon foundational literacy skills, emphasizing phonics, fluent reading of grade-level texts, comprehension strategies, and writing complete sentences. Students engage with literature and informational texts, developing vocabulary and communication skills necessary for academic and social success.

\section{Course Goals}

Upon completion of this course, students will be able to:

\begin{itemize}
    \item Apply phonics skills to decode words
    \item Read grade-level texts with accuracy and fluency
    \item Ask and answer questions about key details in texts
    \item Identify main topics and retell key details of texts
    \item Write opinion pieces, informative/explanatory texts, and narratives
    \item Use conventional spelling for words with common patterns
    \item Participate in collaborative conversations with peers and adults
    \item Produce and expand complete sentences
\end{itemize}

\section{Units of Study}

\subsection{Unit 1: Phonics and Word Study (Ongoing)}

\textbf{Duration:} Integrated throughout the year

\textbf{Topics:}
\begin{itemize}
    \item Short vowel sounds (CVC words)
    \item Long vowel sounds (CVCe words, vowel teams)
    \item Consonant blends and digraphs (bl, cr, sh, th, ch, wh, ph)
    \item R-controlled vowels (ar, er, ir, or, ur)
    \item Common word families (-at, -an, -it, -op, -un)
    \item High-frequency sight words (Dolch and Fry lists)
    \item Decoding two-syllable words
\end{itemize}

\subsection{Unit 2: Reading Foundation Skills (10-12 weeks)}

\textbf{Topics:}
\begin{itemize}
    \item Reading grade-level texts with purpose and understanding
    \item Using decoding strategies automatically
    \item Self-monitoring and self-correcting reading errors
    \item Reading with appropriate expression and pacing
    \item Recognizing and reading irregularly spelled words
    \item Building reading stamina and fluency
\end{itemize}

\subsection{Unit 3: Reading Literature (10-12 weeks)}

\textbf{Topics:}
\begin{itemize}
    \item Asking and answering who, what, where, when, why, and how questions
    \item Retelling stories with key details (beginning, middle, end)
    \item Describing characters, settings, and major events using key details
    \item Identifying words and phrases that suggest feelings or appeal to senses
    \item Comparing and contrasting characters and events in stories
    \item Making predictions and inferences
    \item Understanding story structure (problem and solution)
\end{itemize}

\subsection{Unit 4: Reading Informational Text (8-10 weeks)}

\textbf{Topics:}
\begin{itemize}
    \item Asking and answering questions about key details
    \item Identifying the main topic and retelling key details
    \item Describing connections between individuals, events, or ideas
    \item Using text features (headings, tables of contents, glossaries, icons)
    \item Distinguishing between information provided by pictures and words
    \item Identifying reasons an author gives to support points
    \item Comparing and contrasting two texts on the same topic
\end{itemize}

\subsection{Unit 5: Writing (Ongoing)}

\textbf{Duration:} Integrated throughout the year

\textbf{Topics:}
\begin{itemize}
    \item Writing opinion pieces stating an opinion and supplying a reason
    \item Writing informative/explanatory texts naming a topic and supplying facts
    \item Writing narratives recounting events in sequence with details
    \item Focusing on a topic and strengthening writing with guidance
    \item Using digital tools to produce and publish writing with support
    \item Participating in shared research and writing projects
    \item Recalling information to answer questions
\end{itemize}

\subsection{Unit 6: Language and Grammar (Ongoing)}

\textbf{Duration:} Integrated throughout the year

\textbf{Topics:}
\begin{itemize}
    \item Printing all upper and lowercase letters legibly
    \item Using common, proper, and possessive nouns
    \item Using singular and plural nouns with matching verbs
    \item Using personal, possessive, and indefinite pronouns
    \item Using verbs to convey sense of past, present, and future
    \item Using adjectives and conjunctions
    \item Producing and expanding complete simple and compound sentences
    \item Capitalizing dates and names of people
    \item Using end punctuation and commas in dates and series
    \item Using conventional spelling for words with common patterns
    \item Spelling untaught words phonetically
\end{itemize}

\section{Required Materials}

\begin{itemize}
    \item Leveled readers (Guided Reading Levels D--J)
    \item Picture books and chapter books for read-alouds
    \item Writing journals and graphic organizers
    \item Phonics manipulatives (letter tiles, word sorts)
    \item High-frequency word cards
    \item Classroom library with diverse texts
    \item Digital resources for reading and writing
\end{itemize}

\section{Assessment and Grading}

\begin{center}
\begin{tabular}{|l|c|}
\hline
\textbf{Assessment Category} & \textbf{Percentage} \\
\hline
Reading Assessments (running records, comprehension) & 30\% \\
\hline
Writing Samples and Projects & 30\% \\
\hline
Phonics and Spelling Assessments & 20\% \\
\hline
Class Participation and Discussion & 10\% \\
\hline
Homework and Daily Work & 10\% \\
\hline
\end{tabular}
\end{center}

\chapter{Grade 1 Mathematics}

\section{Course Overview}

\textbf{Grade Level:} Grade 1\\
\textbf{Duration:} Full Year (36 weeks)\\
\textbf{Prerequisites:} Kindergarten or demonstrated readiness

\section{Course Description}

Grade 1 Mathematics focuses on developing number sense, addition and subtraction within 20, place value understanding to 120, measurement, time, and geometry. Students engage in problem-solving, mathematical reasoning, and develop strategies for computational fluency. The curriculum emphasizes conceptual understanding alongside procedural skill.

\section{Course Goals}

Upon completion of this course, students will be able to:

\begin{itemize}
    \item Add and subtract fluently within 10
    \item Solve addition and subtraction word problems within 20
    \item Understand place value for numbers to 120
    \item Add within 100 using concrete models and drawings
    \item Tell and write time to the hour and half-hour
    \item Measure lengths using non-standard and standard units
    \item Organize, represent, and interpret data
    \item Reason with shapes and their attributes
\end{itemize}

\section{Units of Study}

\subsection{Unit 1: Adding, Subtracting, and Working with Data (3-4 weeks)}

\textbf{Topics:}
\begin{itemize}
    \item Adding within 10 using various strategies
    \item Subtracting within 10
    \item Understanding the relationship between addition and subtraction
    \item Representing and interpreting categorical data using charts and graphs
    \item Asking and answering questions about data
\end{itemize}

\subsection{Unit 2: Addition and Subtraction Story Problems (4-5 weeks)}

\textbf{Topics:}
\begin{itemize}
    \item Solving Add To/Result Unknown problems
    \item Solving Take From/Result Unknown problems
    \item Solving Put Together/Total Unknown problems
    \item Solving Put Together/Addend Unknown problems
    \item Using drawings and equations to represent problems
    \item Understanding the meaning of the equal sign
\end{itemize}

\subsection{Unit 3: Adding and Subtracting within 20 (5-6 weeks)}

\textbf{Topics:}
\begin{itemize}
    \item Using strategies to add within 20 (counting on, making ten, doubles)
    \item Using strategies to subtract within 20
    \item Relating addition and subtraction (fact families)
    \item Solving Compare/Difference Unknown problems
    \item Building fluency with sums and differences within 10
    \item Writing equations with unknowns in all positions
\end{itemize}

\subsection{Unit 4: Numbers to 99 (4-5 weeks)}

\textbf{Topics:}
\begin{itemize}
    \item Understanding place value (tens and ones)
    \item Counting to 99 by ones and tens
    \item Reading and writing numbers to 99
    \item Comparing two-digit numbers using $<$, $>$, $=$
    \item Understanding that 10 can be thought of as a bundle of ten ones
    \item Composing and decomposing two-digit numbers
\end{itemize}

\subsection{Unit 5: Adding within 100 (4-5 weeks)}

\textbf{Topics:}
\begin{itemize}
    \item Adding tens and ones
    \item Using concrete models and drawings to represent addition
    \item Recording strategies for adding
    \item Adding two-digit numbers with and without regrouping
\end{itemize}

\subsection{Unit 6: Measurement and Data (4-5 weeks)}

\textbf{Topics:}
\begin{itemize}
    \item Ordering objects by length and expressing comparisons
    \item Measuring lengths using non-standard units
    \item Measuring lengths using standard units (inches, centimeters)
    \item Telling time to the hour and half-hour
    \item Telling and writing time in hours and half-hours
\end{itemize}

\subsection{Unit 7: Geometry and Spatial Reasoning (5-6 weeks)}

\textbf{Topics:}
\begin{itemize}
    \item Identifying and naming two-dimensional shapes
    \item Identifying and naming three-dimensional shapes
    \item Distinguishing between two-dimensional and three-dimensional shapes
    \item Analyzing and comparing shapes
    \item Composing new shapes from existing shapes
    \item Partitioning circles and rectangles into two and four equal shares
\end{itemize}

\subsection{Unit 8: More Operations (3-4 weeks)}

\textbf{Topics:}
\begin{itemize}
    \item Understanding subtraction as taking away
    \item Solving subtraction word problems
    \item Relating addition and subtraction
    \item Building subtraction fact fluency
\end{itemize}

\section{Required Materials}

\begin{itemize}
    \item Manipulatives: base-ten blocks, unit cubes, snap cubes
    \item Number lines and hundreds charts
    \item Place value mats and ten-frames
    \item Measurement tools (rulers, measuring tapes)
    \item Clocks (demonstration and student-sized)
    \item Geometric shapes (2D and 3D)
    \item Math journals
\end{itemize}

\section{Assessment and Grading}

\begin{center}
\begin{tabular}{|l|c|}
\hline
\textbf{Assessment Category} & \textbf{Percentage} \\
\hline
Computation Assessments (fact fluency, operation skills) & 40\% \\
\hline
Problem-Solving and Reasoning & 30\% \\
\hline
Performance Tasks & 15\% \\
\hline
Class Participation and Discussions & 10\% \\
\hline
Homework and Daily Work & 5\% \\
\hline
\end{tabular}
\end{center}

\chapter{Grade 1 Science}

\section{Course Overview}

Grade Level: Grade 1

Duration: Full Year (36 weeks)

Prerequisites: Kindergarten or demonstrated readiness

\section{Course Description}

Grade 1 Science builds on foundational inquiry skills through structured hands-on investigations. Students observe, describe, and classify objects and organisms while developing an understanding of life science, physical science, and earth science concepts appropriate for the grade level.

\section{Course Goals}

Upon completion of this course, students will be able to:

\begin{itemize}
\item Ask questions and make predictions based on observations
\item Plan and conduct simple investigations
\item Record and organize data using drawings and simple charts
\item Communicate findings through oral and written explanations
\item Understand basic needs of living things
\item Investigate properties and states of matter
\item Observe patterns in the sky
\item Understand how the sun, moon, and stars change position
\item Explore sound and light
\end{itemize}

\section{Units of Study}

\subsection{Unit 1: Science Skills and Inquiry (Ongoing)}

Duration: Integrated throughout the year

Topics:
\begin{itemize}
\item Asking questions and making predictions
\item Using tools safely (hand lenses, thermometers, rulers)
\item Observing and describing using all five senses
\item Recording observations through drawings and simple charts
\item Sorting and classifying objects by properties
\item Communicating findings to peers
\end{itemize}

\subsection{Unit 2: Plants and Their Needs (6-7 weeks)}

Topics:
\begin{itemize}
\item Identifying parts of plants (roots, stems, leaves, flowers)
\item Understanding that plants need water, air, light, and nutrients
\item Observing plant growth and development
\item Comparing different types of plants
\item Understanding that plants make their own food
\item Exploring how plants respond to their environment
\end{itemize}

Investigations:
\begin{itemize}
\item Growing plants under different conditions
\item Observing seed germination
\item Comparing plant parts and functions
\item Testing what plants need to grow
\end{itemize}

\subsection{Unit 3: Animals and Their Needs (6-7 weeks)}

Topics:
\begin{itemize}
\item Understanding that animals need food, water, air, and shelter
\item Comparing different animal body parts and their functions
\item Observing animal behaviors
\item Understanding that animals have offspring
\item Classifying animals by characteristics
\item Exploring how animals respond to their environment
\end{itemize}

Investigations:
\begin{itemize}
\item Observing animals and their behaviors
\item Comparing animal body coverings
\item Studying how animals move
\item Creating simple food chains
\end{itemize}

\subsection{Unit 4: Properties of Matter (5-6 weeks)}

Topics:
\begin{itemize}
\item Identifying and describing properties of matter (color, size, shape, texture, weight)
\item Exploring solid, liquid, and gas states of matter
\item Sorting and classifying materials by properties
\item Investigating how materials can change
\item Understanding that objects are made of one or more materials
\end{itemize}

Investigations:
\begin{itemize}
\item Testing properties of different materials
\item Exploring changes in matter (melting, freezing, mixing)
\item Sorting objects by observable properties
\item Investigating sinking and floating
\end{itemize}

\subsection{Unit 5: Sound and Light (5-6 weeks)}

Topics:
\begin{itemize}
\item Understanding that sound is caused by vibrations
\item Exploring how sound travels
\item Investigating pitch and volume
\item Understanding that light is needed to see
\item Exploring shadows
\item Investigating how light travels and reflects
\end{itemize}

Investigations:
\begin{itemize}
\item Creating sounds with different objects
\item Making shadow puppets
\item Exploring mirrors and reflection
\item Testing materials that block or allow light
\end{itemize}

\subsection{Unit 6: Objects in the Sky (5-6 weeks)}

Topics:
\begin{itemize}
\item Observing the sun, moon, and stars
\item Understanding that the sun gives us light and heat
\item Observing patterns of the moon's appearance
\item Understanding day and night
\item Recognizing that the sun appears to move across the sky
\item Observing seasonal changes
\end{itemize}

Investigations:
\begin{itemize}
\item Tracking the sun's position during the day
\item Recording moon observations over time
\item Creating day and night models
\item Observing seasonal weather patterns
\end{itemize}

\subsection{Unit 7: Weather and Seasons (4-5 weeks)}

Topics:
\begin{itemize}
\item Observing and recording daily weather
\item Understanding weather changes
\item Describing the four seasons
\item Understanding how weather affects people and animals
\item Recognizing seasonal patterns
\end{itemize}

Investigations:
\begin{itemize}
\item Daily weather observations and graphing
\item Measuring temperature, wind, and precipitation
\item Comparing seasonal changes
\item Creating weather instruments
\end{itemize}

\section{Required Materials}

\begin{itemize}
\item Science journals or observation notebooks
\item Hand lenses and magnifying glasses
\item Thermometers
\item Rulers and measuring tools
\item Plant seeds and growing containers
\item Balance scales
\item Magnets
\item Mirrors
\item Natural materials (rocks, leaves, sticks)
\item Simple machines (pulleys, levers, inclined planes)
\item Safety equipment
\end{itemize}

\section{Assessment and Grading}

\begin{center}
\begin{tabular}{|l|l|}
\hline
\textbf{Assessment Type} & \textbf{Description} \\
\hline
Science Journals & Drawings and observations documenting learning \\
\hline
Performance Tasks & Hands-on investigations and demonstrations \\
\hline
Oral Explanations & Verbal communication of scientific understanding \\
\hline
Observational Notes & Teacher observations during investigations \\
\hline
Simple Projects & Projects demonstrating conceptual understanding \\
\hline
\end{tabular}
\end{center}

\chapter{Grade 1 Social Studies}

\section{Course Overview}

Grade Level: Grade 1

Duration: Full Year (36 weeks)

Prerequisites: Kindergarten or demonstrated readiness

\section{Course Description}

Grade 1 Social Studies expands students' understanding of family, school, and community while introducing basic geographic concepts, civic responsibility, and economic principles. Students explore their roles as citizens and learn about diverse cultures and traditions.

\section{Course Goals}

Upon completion of this course, students will be able to:

\begin{itemize}
\item Understand family structures and responsibilities
\item Recognize the importance of rules and laws
\item Identify community helpers and their roles
\item Use simple maps and globes
\item Understand basic economic concepts
\item Recognize national symbols and holidays
\item Appreciate cultural diversity
\item Demonstrate good citizenship
\end{itemize}

\section{Units of Study}

\subsection{Unit 1: Our Families (4-5 weeks)}

Topics:
\begin{itemize}
\item Understanding different family structures
\item Identifying family members and their roles
\item Recognizing family responsibilities
\item Learning about family traditions and customs
\item Understanding how families change over time
\item Comparing similarities and differences among families
\end{itemize}

\subsection{Unit 2: Our School Community (4-5 weeks)}

Topics:
\begin{itemize}
\item Understanding school rules and why they matter
\item Identifying school staff and their responsibilities
\item Learning about being a good classmate
\item Understanding school traditions
\item Respecting diversity in our school
\item Developing school pride and spirit
\end{itemize}

\subsection{Unit 3: Our Neighborhood and Community (5-6 weeks)}

Topics:
\begin{itemize}
\item Defining what a community is
\item Identifying places in our community
\item Understanding community rules and laws
\item Learning about community helpers (police, firefighters, doctors, etc.)
\item Understanding how people work together
\item Recognizing community problems and solutions
\end{itemize}

\subsection{Unit 4: Basic Geography and Maps (5-6 weeks)}

Topics:
\begin{itemize}
\item Using simple maps of the classroom and school
\item Understanding map symbols and keys
\item Learning directional words (left, right, near, far)
\item Identifying landforms and bodies of water
\item Using globes to find land and water
\item Understanding that maps represent real places
\end{itemize}

\subsection{Unit 5: Goods, Services, and Economics (4-5 weeks)}

Topics:
\begin{itemize}
\item Understanding the difference between goods and services
\item Learning about jobs and work
\item Understanding wants and needs
\item Recognizing that people earn money by working
\item Making choices about spending
\item Understanding basic concepts of buying and selling
\end{itemize}

\subsection{Unit 6: American Symbols and Holidays (4-5 weeks)}

Topics:
\begin{itemize}
\item Identifying the American flag and what it represents
\item Learning about the Pledge of Allegiance
\item Recognizing other American symbols (Statue of Liberty, Eagle, etc.)
\item Understanding patriotic holidays
\item Learning about important American figures
\item Showing respect for symbols and traditions
\end{itemize}

\subsection{Unit 7: Cultures and Diversity (4-5 weeks)}

Topics:
\begin{itemize}
\item Learning about different cultures and traditions
\item Recognizing cultural celebrations and holidays
\item Understanding that families celebrate in different ways
\item Appreciating diversity and similarities
\item Sharing cultural traditions and foods
\item Respecting differences
\end{itemize}

\subsection{Unit 8: Past, Present, and Future (3-4 weeks)}

Topics:
\begin{itemize}
\item Understanding the concepts of past, present, and future
\item Comparing how things were, are, and will be
\item Learning about changes over time
\item Creating personal timelines
\item Understanding family history
\item Recognizing historical changes in the community
\end{itemize}

\subsection{Unit 9: Being a Good Citizen (Ongoing)}

Duration: Integrated throughout the year

Topics:
\begin{itemize}
\item Understanding rights and responsibilities
\item Following rules and explaining why rules are important
\item Showing respect for others
\item Working cooperatively
\item Helping others and the community
\item Making good choices
\item Resolving conflicts peacefully
\item Taking care of our environment
\end{itemize}

\section{Required Materials}

\begin{itemize}
\item Simple maps (classroom, school, neighborhood)
\item Globe and world map
\item American flag and pictures of national symbols
\item Picture books about families, communities, and cultures
\item Community helper props
\item Photos of historical and current community
\item Cultural artifacts and items
\item Chart paper for class projects
\end{itemize}

\section{Assessment and Grading}

\begin{center}
\begin{tabular}{|l|l|}
\hline
\textbf{Assessment Type} & \textbf{Description} \\
\hline
Observations & Daily observations of citizenship behaviors \\
\hline
Discussions & Class conversations and sharing \\
\hline
Projects & Simple projects (family tree, community map) \\
\hline
Drawings and Writing & Visual and written representations \\
\hline
Role-Playing & Dramatic activities demonstrating understanding \\
\hline
\end{tabular}
\end{center}


\clearpage

\part{Grade 2}

\chapter{Grade 2 English/Language Arts}

\section{Course Overview}

\textbf{Grade Level:} Grade 2\\
\textbf{Duration:} Full Year (36 weeks)\\
\textbf{Prerequisites:} Grade 1 or demonstrated readiness

\section{Course Description}

Grade 2 English/Language Arts strengthens reading fluency and comprehension, writing skills, and language development. Students read and comprehend literature and informational texts with increasing independence, write for multiple purposes, and develop more sophisticated communication skills.

\section{Course Goals}

Upon completion of this course, students will be able to:

\begin{itemize}
    \item Read grade-level texts fluently and with comprehension
    \item Ask and answer literal and inferential questions about texts
    \item Write opinion pieces, informative/explanatory texts, and narratives
    \item Use grade-appropriate grammar and mechanics
    \item Engage in collaborative conversations and presentations
    \item Determine the meaning of unknown words and phrases
    \item Understand how illustrations support text understanding
\end{itemize}

\section{Units of Study}

\subsection{Unit 1: Reading Foundation Skills (Ongoing)}

\textbf{Duration:} Integrated throughout the year

\textbf{Topics:}
\begin{itemize}
    \item Reading with accuracy, fluency, and comprehension
    \item Decoding unfamiliar words
    \item Using phonics and word analysis skills
    \item Recognizing high-frequency words
    \item Self-correcting while reading
\end{itemize}

\subsection{Unit 2: Reading Literature (10-12 weeks)}

\textbf{Topics:}
\begin{itemize}
    \item Asking and answering ``who, what, where, when, why, how'' questions
    \item Retelling main events with key details
    \item Describing characters using key details
    \item Identifying the setting of a story
    \item Identifying the main problem and solution in a story
    \item Describing how characters respond to events
    \item Comparing and contrasting characters and events
\end{itemize}

\subsection{Unit 3: Reading Informational Text (8-10 weeks)}

\textbf{Topics:}
\begin{itemize}
    \item Asking and answering questions about key details
    \item Identifying the main topic of a text
    \item Describing the connection between events
    \item Using text features (headings, table of contents, glossaries)
    \item Distinguishing between information provided by pictures and words
    \item Identifying the purpose of a text
    \item Comparing and contrasting two texts on the same topic
\end{itemize}

\subsection{Unit 4: Writing (Ongoing)}

\textbf{Duration:} Integrated throughout the year

\textbf{Topics:}
\begin{itemize}
    \item Writing opinion pieces with supporting details
    \item Writing informative/explanatory texts with facts and details
    \item Writing narratives with a clear sequence of events
    \item Focusing on a topic and developing it with details
    \item Using transitions to connect ideas
    \item Revising and editing writing
    \item Using digital tools to produce and publish writing
\end{itemize}

\subsection{Unit 5: Language and Grammar (Ongoing)}

\textbf{Duration:} Integrated throughout the year

\textbf{Topics:}
\begin{itemize}
    \item Using nouns, verbs, adjectives, and adverbs appropriately
    \item Using pronouns and antecedents correctly
    \item Using verbs in different tenses
    \item Using coordinating conjunctions (and, or, but)
    \item Producing complete sentences
    \item Using correct capitalization and end punctuation
    \item Using commas in greetings and closings
    \item Using conventional spelling
    \item Determining meanings of unknown words
\end{itemize}

\subsection{Unit 6: Vocabulary Acquisition (Ongoing)}

\textbf{Duration:} Integrated throughout the year

\textbf{Topics:}
\begin{itemize}
    \item Determining the meaning of unknown words using context
    \item Using word relationships and nuances
    \item Understanding word categories and relationships
    \item Using prefixes and suffixes (un-, re-, -ly, -less)
    \item Exploring shades of meaning in synonyms
\end{itemize}

\section{Required Materials}

\begin{itemize}
    \item Leveled readers (Guided Reading Levels J--M)
    \item Picture books and beginning chapter books
    \item Writing folders and graphic organizers
    \item Word study materials
    \item High-frequency word cards
    \item Classroom library with diverse texts
    \item Digital resources for writing and word study
\end{itemize}

\section{Assessment and Grading}

\begin{center}
\begin{tabular}{|l|c|}
\hline
\textbf{Assessment Category} & \textbf{Percentage} \\
\hline
Reading Comprehension & 30\% \\
\hline
Writing Samples and Projects & 30\% \\
\hline
Word Study and Phonics & 15\% \\
\hline
Grammar and Language & 10\% \\
\hline
Class Participation and Discussion & 10\% \\
\hline
Homework and Daily Work & 5\% \\
\hline
\end{tabular}
\end{center}

\chapter{Grade 2 Mathematics}

\section{Course Overview}

\textbf{Grade Level:} Grade 2\\
\textbf{Duration:} Full Year (36 weeks)\\
\textbf{Prerequisites:} Grade 1 or demonstrated readiness

\section{Course Description}

Grade 2 Mathematics builds on number sense and develops fluency with addition and subtraction, introduces multiplication and division concepts, develops understanding of place value to 1000, and explores measurement, time, and data. Students reason about and solve real-world problems using increasingly sophisticated strategies.

\section{Course Goals}

Upon completion of this course, students will be able to:

\begin{itemize}
    \item Fluently add and subtract within 20
    \item Understand place value for numbers to 1000
    \item Add and subtract within 1000 using strategies and algorithms
    \item Work with equal groups of objects to gain foundations for multiplication
    \item Measure and estimate lengths in standard units
    \item Relate addition and subtraction to length
    \item Tell and write time to the nearest five minutes
    \item Solve problems involving money
    \item Represent and interpret data
\end{itemize}

\section{Units of Study}

\subsection{Unit 1: Adding, Subtracting, and Working with Data (3-4 weeks)}

\textbf{Topics:}
\begin{itemize}
    \item Fluently adding and subtracting within 20
    \item Creating and solving story problems
    \item Organizing data in picture graphs and bar graphs
    \item Interpreting data displays
\end{itemize}

\subsection{Unit 2: Addition and Subtraction within 100 (5-6 weeks)}

\textbf{Topics:}
\begin{itemize}
    \item Adding two-digit numbers using strategies (place value, composition, decomposition)
    \item Subtracting two-digit numbers using strategies
    \item Recording strategies with drawings and symbols
    \item Solving real-world addition and subtraction problems
\end{itemize}

\subsection{Unit 3: Measuring Length (4-5 weeks)}

\textbf{Topics:}
\begin{itemize}
    \item Measuring lengths using standard units (inches, feet, centimeters, meters)
    \item Comparing and ordering measurements
    \item Estimating lengths
    \item Measuring and comparing lengths to solve problems
    \item Representing measurement data in line plots
\end{itemize}

\subsection{Unit 4: Addition and Subtraction on the Number Line (3-4 weeks)}

\textbf{Topics:}
\begin{itemize}
    \item Using number lines to model addition
    \item Using number lines to model subtraction
    \item Solving problems using number line representations
\end{itemize}

\subsection{Unit 5: Numbers to 1,000 (4-5 weeks)}

\textbf{Topics:}
\begin{itemize}
    \item Understanding place value (hundreds, tens, ones)
    \item Reading and writing numbers to 1000
    \item Comparing three-digit numbers
    \item Counting by 5s, 10s, and 100s
    \item Composing and decomposing three-digit numbers
\end{itemize}

\subsection{Unit 6: Geometry, Time, and Money (5-6 weeks)}

\textbf{Topics:}
\begin{itemize}
    \item Identifying and describing two-dimensional shapes
    \item Identifying and describing three-dimensional shapes
    \item Partitioning shapes into equal parts
    \item Telling time to the nearest five minutes
    \item Counting money (pennies, nickels, dimes, quarters)
    \item Solving problems involving money
\end{itemize}

\subsection{Unit 7: Adding and Subtracting within 1,000 (5-6 weeks)}

\textbf{Topics:}
\begin{itemize}
    \item Adding three-digit numbers using strategies and algorithms
    \item Subtracting three-digit numbers using strategies and algorithms
    \item Recording strategies
    \item Solving real-world problems
    \item Checking reasonableness of answers
\end{itemize}

\subsection{Unit 8: Equal Groups (4-5 weeks)}

\textbf{Topics:}
\begin{itemize}
    \item Identifying equal groups
    \item Counting equal groups using repeated addition
    \item Using arrays to understand equal groups
    \item Foundations for multiplication and division
\end{itemize}

\subsection{Unit 9: Putting It All Together (3 weeks)}

\textbf{Topics:}
\begin{itemize}
    \item Reviewing and reinforcing key concepts
    \item Solving multi-step problems
    \item Applying strategies to new situations
    \item Problem-solving with real-world contexts
\end{itemize}

\section{Required Materials}

\begin{itemize}
    \item Base-ten blocks and other manipulatives
    \item Number lines and hundreds charts
    \item Place value mats
    \item Rulers, measuring tapes, and meter sticks
    \item Clocks (demonstration and student-sized)
    \item Money manipulatives (coins and bills)
    \item Geometric shapes (2D and 3D)
    \item Graph paper
    \item Math journals
\end{itemize}

\section{Assessment and Grading}

\begin{center}
\begin{tabular}{|l|c|}
\hline
\textbf{Assessment Category} & \textbf{Percentage} \\
\hline
Computation and Number Sense & 40\% \\
\hline
Problem-Solving and Reasoning & 30\% \\
\hline
Performance Tasks & 15\% \\
\hline
Class Participation and Discussions & 10\% \\
\hline
Homework and Daily Work & 5\% \\
\hline
\end{tabular}
\end{center}

\chapter{Grade 2 Science}

\section{Course Overview}

\textbf{Grade Level:} Grade 2\\
\textbf{Duration:} Full Year (36 weeks)\\
\textbf{Prerequisites:} Grade 1 or demonstrated readiness

\section{Course Description}

Grade 2 Science expands inquiry and observation skills through investigations of life science, physical science, and earth and space science. Students plan and conduct investigations, make observations, and develop explanations based on evidence. The curriculum fosters scientific thinking and curiosity about the natural world.

\section{Course Goals}

Upon completion of this course, students will be able to:

\begin{itemize}
    \item Plan and conduct investigations with guidance
    \item Make detailed observations and record data
    \item Construct explanations based on evidence
    \item Communicate scientific findings
    \item Understand plant and animal life cycles
    \item Investigate properties of materials and changes in matter
    \item Observe and record weather patterns and earth's surface
    \item Understand organisms and their relationships with environments
\end{itemize}

\section{Units of Study}

\subsection{Unit 1: Life Cycles (8-9 weeks)}

\textbf{Topics:}
\begin{itemize}
    \item Observing plant growth and development
    \item Understanding plant needs (water, sunlight, soil)
    \item Observing animal behavior and habitats
    \item Understanding animal needs
    \item Comparing life cycles of different organisms
    \item Recording observations over time
\end{itemize}

\textbf{Investigations:}
\begin{itemize}
    \item Growing plants and observing growth
    \item Observing germination and growth rates
    \item Studying animal habitats and behaviors
    \item Creating and maintaining animal habitats
\end{itemize}

\subsection{Unit 2: Properties and Changes in Matter (8-9 weeks)}

\textbf{Topics:}
\begin{itemize}
    \item Observing and describing properties of materials
    \item Comparing different materials
    \item Investigating states of matter (solid, liquid, gas)
    \item Exploring physical and chemical changes
    \item Understanding that materials can be used in different ways
    \item Investigating dissolving and mixing
\end{itemize}

\textbf{Investigations:}
\begin{itemize}
    \item Sorting materials by properties
    \item Testing materials for specific properties
    \item Melting and freezing observations
    \item Dissolving and mixing experiments
    \item Heating and cooling investigations
\end{itemize}

\subsection{Unit 3: Weather and Earth's Surface (8-9 weeks)}

\textbf{Topics:}
\begin{itemize}
    \item Observing daily and seasonal weather patterns
    \item Recording weather data
    \item Understanding weather effects on plants and animals
    \item Investigating Earth's surface (rocks, soil, water)
    \item Understanding the water cycle
    \item Exploring the effects of erosion
\end{itemize}

\textbf{Investigations:}
\begin{itemize}
    \item Daily weather observations and data collection
    \item Creating weather graphs and charts
    \item Examining soil and rock samples
    \item Water cycle demonstrations
    \item Investigating erosion and weathering
\end{itemize}

\subsection{Unit 4: Organisms and Environments (5-6 weeks)}

\textbf{Topics:}
\begin{itemize}
    \item Understanding habitats and environments
    \item Identifying organisms and their adaptations
    \item Understanding food chains and food webs
    \item Recognizing relationships between organisms
    \item Understanding how organisms meet their needs
    \item Exploring human impacts on environments
\end{itemize}

\textbf{Investigations:}
\begin{itemize}
    \item Exploring local ecosystems
    \item Observing food chains
    \item Studying organism adaptations
    \item Environmental impact investigations
\end{itemize}

\section{Required Materials}

\begin{itemize}
    \item Science journals and observation notebooks
    \item Hand lenses and magnifying glasses
    \item Thermometers
    \item Balance scales
    \item Measuring cups and tools
    \item Plant seeds and growing containers
    \item Weather observation charts
    \item Natural materials (rocks, soil, leaves, water)
    \item Magnets
    \item Art supplies for drawings
    \item Safety equipment
\end{itemize}

\section{Assessment and Grading}

\begin{center}
\begin{tabular}{|l|c|}
\hline
\textbf{Assessment Category} & \textbf{Percentage} \\
\hline
Scientific Investigations and Data Collection & 35\% \\
\hline
Science Journals and Observations & 25\% \\
\hline
Explanations and Reasoning & 20\% \\
\hline
Class Participation and Collaboration & 15\% \\
\hline
Homework and Daily Work & 5\% \\
\hline
\end{tabular}
\end{center}

\chapter{Grade 2 Social Studies}

\section{Course Overview}

\textbf{Grade Level:} Grade 2\\
\textbf{Duration:} Full Year (36 weeks)\\
\textbf{Prerequisites:} Grade 1 or demonstrated readiness

\section{Course Description}

Grade 2 Social Studies expands understanding of communities, neighborhoods, and how people meet their needs. Students explore roles in families and schools, community helpers, neighborhoods, and cultural diversity. The curriculum emphasizes map skills, economics, civics, and appreciation for different peoples and places.

\section{Course Goals}

Upon completion of this course, students will be able to:

\begin{itemize}
    \item Understand their role in family, school, and community
    \item Identify community helpers and their roles
    \item Read and create simple maps
    \item Understand basic economic concepts
    \item Appreciate cultural diversity
    \item Identify historical events and people
    \item Demonstrate citizenship behaviors
\end{itemize}

\section{Units of Study}

\subsection{Unit 1: Families and How They Change (3-4 weeks)}

\textbf{Topics:}
\begin{itemize}
    \item Understanding family structures and roles
    \item Learning about family responsibilities
    \item Understanding how families meet basic needs
    \item Exploring family histories and traditions
    \item Comparing different types of families
\end{itemize}

\subsection{Unit 2: School and Community (4-5 weeks)}

\textbf{Topics:}
\begin{itemize}
    \item Understanding school rules and routines
    \item Identifying school staff and their roles
    \item Learning about classroom responsibilities
    \item Understanding the concept of community
    \item Identifying places in the community
    \item Learning about community resources
\end{itemize}

\subsection{Unit 3: Community Helpers (3-4 weeks)}

\textbf{Topics:}
\begin{itemize}
    \item Identifying community helpers
    \item Understanding what community helpers do
    \item Learning about different jobs in the community
    \item Understanding how community helpers serve others
    \item Exploring community services
\end{itemize}

\subsection{Unit 4: Neighborhoods and Maps (4-5 weeks)}

\textbf{Topics:}
\begin{itemize}
    \item Understanding what a neighborhood is
    \item Using and creating simple maps
    \item Understanding map symbols and keys
    \item Identifying locations on maps
    \item Understanding directions (north, south, east, west)
    \item Exploring neighborhoods in maps
\end{itemize}

\subsection{Unit 5: Goods and Services (4-5 weeks)}

\textbf{Topics:}
\begin{itemize}
    \item Understanding goods and services
    \item Learning about jobs and work
    \item Understanding wants and needs
    \item Exploring basic economics (buying and selling)
    \item Understanding money and its uses
    \item Learning about making choices
\end{itemize}

\subsection{Unit 6: People and Cultures (5-6 weeks)}

\textbf{Topics:}
\begin{itemize}
    \item Learning about different cultures
    \item Exploring traditions and customs
    \item Celebrating cultural diversity
    \item Understanding how people are similar and different
    \item Exploring holidays from different cultures
    \item Appreciating cultural contributions
\end{itemize}

\subsection{Unit 7: Then and Now (3-4 weeks)}

\textbf{Topics:}
\begin{itemize}
    \item Understanding then and now
    \item Comparing life in the past and present
    \item Learning about historical figures
    \item Understanding how communities change
    \item Exploring technology changes
\end{itemize}

\section{Required Materials}

\begin{itemize}
    \item Maps (community, state, country, world)
    \item Globes
    \item Picture books about communities and cultures
    \item Historical photographs and artifacts
    \item Community helper props and costumes
    \item Chart paper for class books
    \item Cultural artifacts and items
\end{itemize}

\section{Assessment and Grading}

\begin{center}
\begin{tabular}{|l|c|}
\hline
\textbf{Assessment Category} & \textbf{Percentage} \\
\hline
Concept Understanding and Application & 30\% \\
\hline
Projects and Presentations & 25\% \\
\hline
Map and Geography Skills & 20\% \\
\hline
Class Participation and Discussions & 15\% \\
\hline
Homework and Daily Work & 10\% \\
\hline
\end{tabular}
\end{center}

\clearpage

\part{Grade 3}

\chapter{Grade 3 English/Language Arts}

\section{Course Overview}

\textbf{Grade Level:} Grade 3\\
\textbf{Duration:} Full Year (36 weeks)\\
\textbf{Prerequisites:} Grade 2 or demonstrated readiness

\section{Course Description}

Grade 3 English/Language Arts emphasizes developing fluent readers who comprehend complex literature and informational texts. Students develop more sophisticated writing skills across multiple genres, grammar knowledge, and communication abilities. The curriculum supports critical thinking, analysis, and creative expression.

\section{Course Goals}

Upon completion of this course, students will be able to:

\begin{itemize}
    \item Read and comprehend grade-level literature and informational texts
    \item Ask and answer both literal and inferential questions
    \item Describe characters, settings, and events in detail
    \item Identify main ideas and supporting details
    \item Write well-developed opinion, informative, and narrative texts
    \item Use proper grammar, capitalization, and punctuation
    \item Engage in collaborative conversations and presentations
    \item Determine meanings of new words and phrases
\end{itemize}

\section{Units of Study}

\subsection{Unit 1: Reading Literature (10-12 weeks)}

\textbf{Topics:}
\begin{itemize}
    \item Asking and answering literal and inferential questions
    \item Describing characters' traits, motivations, and feelings
    \item Describing setting and how it affects the story
    \item Identifying plot elements (problem, solution, theme)
    \item Making predictions and inferences
    \item Understanding author's purpose and tone
    \item Comparing and contrasting characters and stories
    \item Identifying literary elements and techniques
\end{itemize}

\subsection{Unit 2: Reading Informational Text (8-10 weeks)}

\textbf{Topics:}
\begin{itemize}
    \item Asking and answering literal and inferential questions
    \item Identifying main topic and supporting details
    \item Describing the sequence of events or steps
    \item Using text features to locate information
    \item Understanding cause and effect relationships
    \item Comparing and contrasting information
    \item Understanding author's purpose
    \item Evaluating information and drawing conclusions
\end{itemize}

\subsection{Unit 3: Writing (Ongoing)}

\textbf{Duration:} Integrated throughout the year

\textbf{Topics:}
\begin{itemize}
    \item Writing opinion pieces with organized reasoning
    \item Writing informative/explanatory texts with facts and examples
    \item Writing narrative texts with descriptive details
    \item Organizing writing with clear structure
    \item Using transitions to connect ideas
    \item Providing details and examples to support writing
    \item Revising and editing for clarity and effectiveness
    \item Publishing and sharing writing
\end{itemize}

\subsection{Unit 4: Language and Grammar (Ongoing)}

\textbf{Duration:} Integrated throughout the year

\textbf{Topics:}
\begin{itemize}
    \item Using nouns, pronouns, verbs, adjectives, and adverbs correctly
    \item Using subject-verb agreement
    \item Using verb tenses appropriately
    \item Using conjunctions and subordinating conjunctions
    \item Using commas appropriately
    \item Using quotation marks in dialogue
    \item Using capitalization correctly
    \item Using and spelling words correctly
\end{itemize}

\subsection{Unit 5: Vocabulary and Word Study (Ongoing)}

\textbf{Duration:} Integrated throughout the year

\textbf{Topics:}
\begin{itemize}
    \item Determining word meanings using context
    \item Understanding word relationships and nuances
    \item Using root words and affixes
    \item Understanding compound words
    \item Exploring synonyms and antonyms
    \item Understanding multiple-meaning words
\end{itemize}

\section{Required Materials}

\begin{itemize}
    \item Leveled readers (Guided Reading Levels M--P)
    \item Chapter books and novels
    \item Picture books and poetry collections
    \item Writing folders and graphic organizers
    \item Word study and phonics materials
    \item Classroom library with diverse texts
    \item Digital resources for reading and writing
\end{itemize}

\section{Assessment and Grading}

\begin{center}
\begin{tabular}{|l|c|}
\hline
\textbf{Assessment Category} & \textbf{Percentage} \\
\hline
Reading Comprehension and Analysis & 35\% \\
\hline
Writing Samples and Projects & 35\% \\
\hline
Grammar, Language, and Vocabulary & 15\% \\
\hline
Class Participation and Discussion & 10\% \\
\hline
Homework and Daily Work & 5\% \\
\hline
\end{tabular}
\end{center}

\chapter{Grade 3 Mathematics}

\section{Course Overview}

\textbf{Grade Level:} Grade 3\\
\textbf{Duration:} Full Year (36 weeks)\\
\textbf{Prerequisites:} Grade 2 or demonstrated readiness

\section{Course Description}

Grade 3 Mathematics develops deeper understanding of multiplication and division, represents and interprets data, understands properties of multiplication, solves multiplication and division word problems, and develops computational strategies. The curriculum emphasizes conceptual understanding, procedural fluency, and application to real-world situations.

\section{Course Goals}

Upon completion of this course, students will be able to:

\begin{itemize}
    \item Develop fluency with multiplication and division facts
    \item Understand place value for numbers up to 1000
    \item Fluently add and subtract within 1000
    \item Represent and solve multiplication and division problems
    \item Understand properties of multiplication and the relationship between multiplication and division
    \item Multiply and divide within 100
    \item Solve problems involving measurement and data
    \item Develop understanding of fractions
    \item Reason with shapes and their attributes
\end{itemize}

\section{Units of Study}

\subsection{Unit 1: Introducing Multiplication (4-5 weeks)}

\textbf{Topics:}
\begin{itemize}
    \item Understanding repeated addition as multiplication
    \item Understanding rows and columns
    \item Using arrays to represent multiplication
    \item Understanding commutative property
    \item Writing multiplication equations
    \item Connecting repeated addition to multiplication
\end{itemize}

\subsection{Unit 2: Area and Multiplication (4-5 weeks)}

\textbf{Topics:}
\begin{itemize}
    \item Measuring area in square units
    \item Using arrays to find area
    \item Relating multiplication to area
    \item Solving real-world area problems
    \item Representing area with multiplication equations
\end{itemize}

\subsection{Unit 3: Wrapping Up Addition and Subtraction within 1,000 (4-5 weeks)}

\textbf{Topics:}
\begin{itemize}
    \item Using efficient strategies for addition and subtraction
    \item Recording strategies with diagrams and equations
    \item Solving multi-digit addition and subtraction problems
    \item Estimating sums and differences
    \item Checking reasonableness of answers
\end{itemize}

\subsection{Unit 4: Relating Multiplication to Division (5-6 weeks)}

\textbf{Topics:}
\begin{itemize}
    \item Understanding division as equal sharing
    \item Understanding division as partitioning
    \item Understanding the relationship between multiplication and division
    \item Using arrays to represent division
    \item Writing division equations
    \item Solving division word problems
\end{itemize}

\subsection{Unit 5: Fractions as Numbers (6-7 weeks)}

\textbf{Topics:}
\begin{itemize}
    \item Understanding unit fractions
    \item Representing fractions with models and drawings
    \item Understanding fractions on a number line
    \item Comparing fractions with the same numerator or denominator
    \item Understanding equivalent fractions
    \item Using fractions in real-world contexts
\end{itemize}

\subsection{Unit 6: Measuring Length, Time, Liquid Volume, and Weight (5-6 weeks)}

\textbf{Topics:}
\begin{itemize}
    \item Measuring lengths using standard units
    \item Estimating and measuring distances
    \item Telling time to the nearest minute
    \item Measuring elapsed time
    \item Measuring liquid volume (capacity)
    \item Measuring and comparing weights
    \item Solving measurement word problems
\end{itemize}

\subsection{Unit 7: Two-dimensional Shapes and Perimeter (4-5 weeks)}

\textbf{Topics:}
\begin{itemize}
    \item Identifying and categorizing shapes
    \item Understanding attributes of shapes
    \item Identifying and drawing quadrilaterals
    \item Understanding perimeter
    \item Measuring perimeter
    \item Solving perimeter problems
    \item Understanding the relationship between area and perimeter
\end{itemize}

\subsection{Unit 8: Putting It All Together (3-4 weeks)}

\textbf{Topics:}
\begin{itemize}
    \item Reviewing and reinforcing key concepts
    \item Solving multi-step word problems
    \item Applying strategies to new situations
    \item Real-world problem-solving
    \item Assessment and reflection
\end{itemize}

\section{Required Materials}

\begin{itemize}
    \item Array cards and multiplication models
    \item Base-ten blocks and manipulatives
    \item Fraction models and tiles
    \item Rulers, measuring tapes, and meter sticks
    \item Clocks (demonstration and student-sized)
    \item Measuring cups and containers
    \item Scales for weighing
    \item Area and perimeter grids
    \item Geometric shapes (2D and 3D)
    \item Graph paper
    \item Math journals
\end{itemize}

\section{Assessment and Grading}

\begin{center}
\begin{tabular}{|l|c|}
\hline
\textbf{Assessment Category} & \textbf{Percentage} \\
\hline
Multiplication, Division, and Computational Fluency & 35\% \\
\hline
Problem-Solving and Reasoning & 30\% \\
\hline
Measurement, Data, and Geometry & 15\% \\
\hline
Class Participation and Discussions & 10\% \\
\hline
Homework and Daily Work & 10\% \\
\hline
\end{tabular}
\end{center}

\chapter{Grade 3 Science}

\section{Course Overview}

\textbf{Grade Level:} Grade 3\\
\textbf{Duration:} Full Year (36 weeks)\\
\textbf{Prerequisites:} Grade 2 or demonstrated readiness

\section{Course Description}

Grade 3 Science develops more sophisticated inquiry skills through structured investigations. Students conduct experiments, collect and organize data, and develop evidence-based explanations. The curriculum covers life science topics including ecosystems and adaptations, physical science topics including forces and energy, and earth science topics.

\section{Course Goals}

Upon completion of this course, students will be able to:

\begin{itemize}
    \item Plan and conduct scientific investigations
    \item Collect, organize, and analyze data
    \item Develop explanations supported by evidence
    \item Communicate scientific findings
    \item Understand plant and animal adaptations
    \item Understand food chains and energy flow
    \item Understand properties of matter and energy transfer
    \item Understand forces and motion
    \item Understand rocks, landforms, and earth processes
\end{itemize}

\section{Units of Study}

The Grade 3 Science curriculum includes investigations in the following areas:

\begin{itemize}
    \item Life Cycles and Adaptations
    \item Food Chains and Ecosystems
    \item Matter and Energy Transfer
    \item Forces and Motion
    \item Rocks and Landforms
\end{itemize}

(Detailed unit plans available in full curriculum document)

\section{Assessment and Grading}

\begin{center}
\begin{tabular}{|l|c|}
\hline
\textbf{Assessment Category} & \textbf{Percentage} \\
\hline
Investigations and Experiments & 40\% \\
\hline
Data Collection and Analysis & 25\% \\
\hline
Scientific Explanations and Reasoning & 20\% \\
\hline
Class Participation and Safety & 10\% \\
\hline
Homework and Journals & 5\% \\
\hline
\end{tabular}
\end{center}

\chapter{Grade 3 Social Studies}

\section{Course Overview}

\textbf{Grade Level:} Grade 3\\
\textbf{Duration:} Full Year (36 weeks)\\
\textbf{Prerequisites:} Grade 2 or demonstrated readiness

\section{Course Description}

Grade 3 Social Studies expands geographic and historical understanding. Students explore communities, regions, and their characteristics. The curriculum includes map skills, cultural diversity, economic activities, government and citizenship, and historical events.

\section{Course Goals}

Upon completion of this course, students will be able to:

\begin{itemize}
    \item Understand characteristics of different communities
    \item Use maps and geographic tools
    \item Understand human and physical geography
    \item Understand basic economic activities
    \item Understand government and civic responsibilities
    \item Understand cultural diversity
    \item Understand historical events and figures
    \item Demonstrate citizenship behaviors
\end{itemize}

\section{Units of Study}

The Grade 3 Social Studies curriculum includes study of:

\begin{itemize}
    \item Community Types and Characteristics
    \item Maps and Geographic Skills
    \item Regions and Regional Characteristics
    \item Economic Activities and Jobs
    \item Government and Civics
    \item Cultural Groups and Diversity
    \item State History and Local History
\end{itemize}

\subsection{Unit 1: Life Cycles and Adaptations (7-8 weeks)}

Topics:
\begin{itemize}
\item Understanding complete and incomplete metamorphosis
\item Observing life cycles of plants and animals
\item Investigating inherited traits
\item Exploring adaptations that help organisms survive
\item Understanding that organisms have different life cycles
\item Comparing life spans of different organisms
\end{itemize}

Investigations:
\begin{itemize}
\item Observing butterfly or frog life cycles
\item Growing plants and documenting life cycle stages
\item Studying adaptations of local organisms
\item Comparing inherited traits in organisms
\end{itemize}

\subsection{Unit 2: Food Chains and Ecosystems (7-8 weeks)}

Topics:
\begin{itemize}
\item Understanding producers, consumers, and decomposers
\item Creating and analyzing food chains
\item Understanding energy flow in ecosystems
\item Exploring different habitats and ecosystems
\item Understanding interdependence of organisms
\item Learning about predator-prey relationships
\end{itemize}

Investigations:
\begin{itemize}
\item Creating food chain models
\item Observing local ecosystems
\item Studying decomposition
\item Building terrariums to observe ecosystems
\end{itemize}

\subsection{Unit 3: Matter and Energy Transfer (7-8 weeks)}

Topics:
\begin{itemize}
\item Investigating properties of matter (mass, volume, density)
\item Exploring states of matter and phase changes
\item Understanding mixtures and solutions
\item Investigating heat energy and temperature
\item Understanding that matter is conserved
\item Exploring physical and chemical changes
\end{itemize}

Investigations:
\begin{itemize}
\item Testing properties of different materials
\item Investigating phase changes (melting, evaporation, condensation)
\item Creating mixtures and solutions
\item Measuring temperature changes
\end{itemize}

\subsection{Unit 4: Forces and Motion (6-7 weeks)}

Topics:
\begin{itemize}
\item Understanding balanced and unbalanced forces
\item Investigating friction, gravity, and magnetism
\item Exploring speed and direction of motion
\item Understanding position and motion
\item Investigating cause and effect relationships
\item Exploring simple machines
\end{itemize}

Investigations:
\begin{itemize}
\item Testing forces with ramps and surfaces
\item Investigating magnetic forces
\item Measuring speed and distance
\item Building and testing simple machines
\end{itemize}

\subsection{Unit 5: Rocks and Landforms (5-6 weeks)}

Topics:
\begin{itemize}
\item Identifying types of rocks (igneous, sedimentary, metamorphic)
\item Understanding the rock cycle
\item Exploring weathering and erosion
\item Investigating soil formation
\item Understanding landforms and how they change
\item Learning about fossils and what they tell us
\end{itemize}

Investigations:
\begin{itemize}
\item Observing and classifying rocks
\item Modeling weathering and erosion
\item Investigating soil composition
\item Creating fossils
\end{itemize}

\section{Required Materials}

\begin{itemize}
\item Science journals
\item Hand lenses and microscopes
\item Thermometers and rulers
\item Balance scales and graduated cylinders
\item Rock and mineral samples
\item Organism specimens or pictures
\item Magnet sets
\item Ramps and force measurement tools
\item Soil samples
\item Safety equipment
\end{itemize}


\section{Assessment and Grading}

\begin{center}
\begin{tabular}{|l|c|}
\hline
\textbf{Assessment Category} & \textbf{Percentage} \\
\hline
Concept Understanding & 30\% \\
\hline
Projects and Presentations & 25\% \\
\hline
Map and Geographic Skills & 20\% \\
\hline
Class Participation and Discussions & 15\% \\
\hline
Homework and Daily Work & 10\% \\
\hline
\end{tabular}
\end{center}

\clearpage

\part{Grade 4}

\chapter{Grade 4 English/Language Arts}

\section{Course Overview}

\textbf{Grade Level:} Grade 4\\
\textbf{Duration:} Full Year (36 weeks)\\
\textbf{Prerequisites:} Grade 3 or demonstrated readiness

\section{Course Description}

Grade 4 English/Language Arts focuses on developing independent readers and writers. Students read diverse literature and informational texts, analyze author's craft, write extended texts in multiple genres, and communicate effectively. The curriculum emphasizes critical thinking, interpretation, and sophisticated communication skills.

\section{Course Goals}

Upon completion of this course, students will be able to:

\begin{itemize}
    \item Read and comprehend grade-level complex texts
    \item Determine main ideas and supporting details
    \item Make inferences and draw conclusions
    \item Analyze author's purpose and point of view
    \item Write well-organized opinion, informative, and narrative texts
    \item Use correct grammar, mechanics, and conventions
    \item Engage in collaborative discussions and presentations
    \item Expand vocabulary and use precise language
\end{itemize}

\section{Units of Study}

(Detailed units covering reading literature, reading informational text, writing, language and grammar)

\section{Assessment and Grading}

\begin{center}
\begin{tabular}{|l|c|}
\hline
\textbf{Assessment Category} & \textbf{Percentage} \\
\hline
Reading Comprehension and Analysis & 35\% \\
\hline
Writing Samples and Projects & 40\% \\
\hline
Grammar, Language, and Vocabulary & 15\% \\
\hline
Class Participation & 10\% \\
\hline
\end{tabular}
\end{center}

\section{Units of Study}

\subsection{Unit 1: Reading Literature and Comprehension Strategies (10-12 weeks)}

Topics:
\begin{itemize}
\item Understanding character traits and motivations
\item Identifying themes and central messages
\item Describing settings in detail
\item Analyzing story elements (plot, conflict, resolution)
\item Making inferences and predictions
\item Comparing and contrasting characters and stories
\item Understanding point of view
\item Reading fluency and expression
\end{itemize}

\subsection{Unit 2: Reading Informational Text (8-10 weeks)}

Topics:
\begin{itemize}
\item Identifying main ideas and supporting details
\item Understanding text structure (cause and effect, sequence, comparison)
\item Using text features (headings, captions, diagrams, glossaries)
\item Describing connections between events, ideas, and information
\item Understanding author's purpose
\item Drawing conclusions from text
\item Using nonfiction text features to locate information
\end{itemize}

\subsection{Unit 3: Writing (Ongoing)}

Duration: Integrated throughout the year

Topics:
\begin{itemize}
\item Writing opinion pieces with supporting reasons
\item Writing informative/explanatory texts with facts and details
\item Writing narratives with description and dialogue
\item Planning, drafting, and revising writing
\item Using transitional words and phrases
\item Using precise language and vocabulary
\item Publishing writing for various audiences
\item Proofreading and editing
\end{itemize}

\subsection{Unit 4: Speaking and Listening (Ongoing)}

Duration: Integrated throughout the year

Topics:
\begin{itemize}
\item Engaging in collaborative discussions
\item Following rules for discussions
\item Building on others' ideas
\item Speaking clearly and at appropriate pace
\item Using multimedia to enhance presentations
\item Asking and answering questions
\item Distinguishing fact from opinion
\end{itemize}

\subsection{Unit 5: Grammar and Language Conventions (Ongoing)}

Duration: Integrated throughout the year

Topics:
\begin{itemize}
\item Using nouns, pronouns, verbs, adjectives, and adverbs correctly
\item Using coordinating and subordinating conjunctions
\item Forming and using verb tenses
\item Capitalizing and punctuating correctly
\item Using commas, quotation marks, and apostrophes
\item Sentence combining and varying sentence structure
\item Proofreading for errors
\end{itemize}

\subsection{Unit 6: Vocabulary and Word Study (Ongoing)}

Duration: Integrated throughout the year

Topics:
\begin{itemize}
\item Using context clues to determine word meanings
\item Understanding word relationships and nuances
\item Using prefixes, suffixes, and root words
\item Sorting and classifying words
\item Understanding synonyms and antonyms
\item Using reference materials (dictionaries, thesauruses)
\item Building academic vocabulary
\item Understanding content-specific vocabulary
\end{itemize}

\chapter{Grade 4 Mathematics}

\section{Course Overview}

\textbf{Grade Level:} Grade 4\\
\textbf{Duration:} Full Year (36 weeks)\\
\textbf{Prerequisites:} Grade 3 or demonstrated readiness

\section{Course Description}

Grade 4 Mathematics develops fluency with multi-digit multiplication and division, extends place value to hundred thousands, develops understanding of fractions and decimals, and applies mathematical thinking to measurement and data problems.

\section{Course Goals}

Upon completion of this course, students will be able to:

\begin{itemize}
    \item Fluently multiply and divide within 100
    \item Solve multi-digit multiplication and division problems
    \item Understand place value through hundred thousands
    \item Add and subtract multi-digit numbers fluently
    \item Understand equivalent fractions and ordering fractions
    \item Understand decimal notation and place value
    \item Measure lengths, convert measurements, and solve problems
    \item Represent and interpret data
    \item Understand angles and geometric properties
\end{itemize}

\section{Units of Study}

(Detailed units covering place value, operations, fractions, decimals, measurement, data, and geometry)

\section{Assessment and Grading}

\begin{center}
\begin{tabular}{|l|c|}
\hline
\textbf{Assessment Category} & \textbf{Percentage} \\
\hline
Computational Skills and Fluency & 35\% \\
\hline
Problem-Solving & 30\% \\
\hline
Measurement, Data, and Geometry & 20\% \\
\hline
Class Participation & 10\% \\
\hline
Daily Work & 5\% \\
\hline
\end{tabular}
\end{center}

\subsection{Unit 1: Multiplication Concepts and Strategies (5-6 weeks)}

Topics:
\begin{itemize}
\item Understanding multiplication as repeated addition
\item Exploring arrays and area models
\item Using skip counting strategies
\item Building multiplication facts (2s, 5s, 10s)
\item Using commutative and associative properties
\item Understanding and applying distributive property
\item Multiplying by multiples of 10
\end{itemize}

\subsection{Unit 2: Division Concepts and Strategies (5-6 weeks)}

Topics:
\begin{itemize}
\item Understanding division as equal sharing and grouping
\item Exploring the relationship between multiplication and division
\item Dividing with remainders
\item Building division facts
\item Interpreting remainders in context
\item Using arrays and area models for division
\item Solving one-step division problems
\end{itemize}

\subsection{Unit 3: Multi-Digit Multiplication (6-7 weeks)}

Topics:
\begin{itemize}
\item Multiplying two-digit by one-digit numbers
\item Multiplying two-digit by two-digit numbers
\item Using arrays and area models
\item Understanding place value in multiplication
\item Estimating products
\item Using strategies for efficient computation
\item Solving multiplication word problems
\end{itemize}

\subsection{Unit 4: Fractions (6-7 weeks)}

Topics:
\begin{itemize}
\item Understanding fractions as parts of a whole
\item Identifying and naming unit fractions
\item Understanding equivalent fractions
\item Comparing and ordering fractions
\item Fractions on number lines
\item Adding and subtracting fractions with like denominators
\item Understanding fractions in real-world contexts
\end{itemize}

\subsection{Unit 5: Decimals and Money (5-6 weeks)}

Topics:
\begin{itemize}
\item Understanding tenths and hundredths
\item Relating fractions to decimals
\item Reading and writing decimal notation
\item Comparing and ordering decimals
\item Adding and subtracting money
\item Solving word problems involving money
\item Making change
\end{itemize}

\subsection{Unit 6: Measurement (5-6 weeks)}

Topics:
\begin{itemize}
\item Measuring length in inches, feet, centimeters, and meters
\item Understanding perimeter and area
\item Finding area of rectangles
\item Measuring weight and mass
\item Measuring liquid volume
\item Reading and using rulers and scales
\item Solving measurement word problems
\end{itemize}

\subsection{Unit 7: Data and Graphing (4-5 weeks)}

Topics:
\begin{itemize}
\item Collecting and organizing data
\item Creating bar graphs and pictographs
\item Creating line plots
\item Reading and interpreting graphs
\item Comparing data sets
\item Understanding mean, median, mode, and range
\item Drawing conclusions from data
\end{itemize}

\subsection{Unit 8: Geometry (5-6 weeks)}

Topics:
\begin{itemize}
\item Identifying points, lines, line segments, and rays
\item Understanding angles and right angles
\item Classifying triangles and quadrilaterals
\item Understanding parallel and perpendicular lines
\item Recognizing symmetry
\item Composing and decomposing shapes
\item Understanding congruence and similarity
\end{itemize}


\chapter{Grade 4 Science}

\section{Course Overview}

\textbf{Grade Level:} Grade 4\\
\textbf{Duration:} Full Year (36 weeks)\\
\textbf{Prerequisites:} Grade 3 or demonstrated readiness

\section{Course Description}

Grade 4 Science focuses on developing deeper understanding of life systems, physical systems, and earth systems. Students plan investigations, collect data systematically, and develop evidence-based explanations of scientific phenomena.

\section{Course Goals}

Upon completion of this course, students will be able to:

\begin{itemize}
    \item Design and conduct scientific investigations
    \item Analyze and interpret data
    \item Develop and support explanations using evidence
    \item Communicate scientific ideas clearly
    \item Understand animal and plant structures and functions
    \item Understand inherited traits and variations
    \item Understand energy transfer and transformations
    \item Understand earth materials and processes
    \item Understand that technology solves human problems
\end{itemize}

\section{Assessment and Grading}

\begin{center}
\begin{tabular}{|l|c|}
\hline
\textbf{Assessment Category} & \textbf{Percentage} \\
\hline
Investigations and Experiments & 40\% \\
\hline
Data Analysis & 25\% \\
\hline
Scientific Reasoning & 20\% \\
\hline
Participation and Safety & 10\% \\
\hline
Science Journals & 5\% \\
\hline
\end{tabular}
\end{center}

\section{Units of Study}

\subsection{Unit 1: Energy and Motion (6-7 weeks)}

Topics:
\begin{itemize}
\item Understanding different forms of energy (light, sound, heat, mechanical)
\item Exploring transfer of energy
\item Understanding motion and speed
\item Investigating forces and their effects
\item Exploring friction, gravity, and magnetism
\item Understanding simple machines (lever, pulley, inclined plane)
\item Investigating cause and effect relationships
\end{itemize}

Investigations:
\begin{itemize}
\item Building simple machines
\item Measuring speed and distance
\item Investigating friction on different surfaces
\item Exploring magnetic forces and fields
\end{itemize}

\subsection{Unit 2: Waves and Sound (5-6 weeks)}

Topics:
\begin{itemize}
\item Understanding that sound is produced by vibration
\item Investigating sound waves and properties
\item Understanding pitch and volume
\item Exploring how sound travels
\item Understanding properties of waves
\item Investigating reflection and absorption of sound
\item Connecting sound to communication
\end{itemize}

Investigations:
\begin{itemize}
\item Creating sounds with different objects
\item Investigating pitch variations
\item Exploring sound in different media
\end{itemize}

\subsection{Unit 3: Weather and Climate (6-7 weeks)}

Topics:
\begin{itemize}
\item Understanding weather patterns and weather prediction
\item Measuring weather variables (temperature, precipitation, wind)
\item Understanding the water cycle
\item Exploring clouds and precipitation
\item Understanding climate and seasons
\item Investigating severe weather
\item Understanding effects of weather on living things
\end{itemize}

Investigations:
\begin{itemize}
\item Daily weather observations and data collection
\item Creating weather instruments
\item Modeling the water cycle
\item Analyzing weather patterns over time
\end{itemize}

\subsection{Unit 4: Rocks, Soil, and Fossils (6-7 weeks)}

Topics:
\begin{itemize}
\item Identifying types of rocks and minerals
\item Understanding the rock cycle
\item Exploring weathering and erosion
\item Understanding soil composition and formation
\item Investigating properties of soil
\item Understanding fossils and fossilization
\item Learning about geological time
\end{itemize}

Investigations:
\begin{itemize}
\item Classifying rocks and minerals
\item Modeling weathering and erosion
\item Investigating soil properties
\item Creating fossil models
\end{itemize}

\subsection{Unit 5: Earth's Features and Systems (5-6 weeks)}

Topics:
\begin{itemize}
\item Understanding landforms and bodies of water
\item Exploring maps and geography
\item Understanding plate tectonics basics
\item Investigating natural resources
\item Understanding conservation of resources
\item Exploring Earth's layers
\end{itemize}

Investigations:
\begin{itemize}
\item Creating topographic models
\item Mapping landforms
\item Investigating natural resources
\end{itemize}

\subsection{Unit 6: Life Cycles and Ecosystems (6-7 weeks)}

Topics:
\begin{itemize}
\item Understanding life cycles of plants and animals
\item Exploring food chains and food webs
\item Understanding energy flow in ecosystems
\item Investigating habitats and adaptations
\item Understanding interdependence of organisms
\item Exploring different ecosystems
\item Understanding ecological niches
\end{itemize}

Investigations:
\begin{itemize}
\item Observing life cycles in nature
\item Creating food webs and chains
\item Building ecosystem models
\item Studying local habitat
\end{itemize}

\subsection{Unit 7: Properties and Changes in Matter (5-6 weeks)}

Topics:
\begin{itemize}
\item Understanding physical and chemical properties
\item Investigating states of matter
\item Exploring phase changes
\item Understanding mixtures and solutions
\item Investigating chemical reactions
\item Understanding conservation of matter
\item Exploring density and buoyancy
\end{itemize}

Investigations:
\begin{itemize}
\item Testing properties of materials
\item Investigating phase changes
\item Creating mixtures and solutions
\item Observing chemical reactions
\end{itemize}

\chapter{Grade 4 Social Studies}

\section{Course Overview}

\textbf{Grade Level:} Grade 4\\
\textbf{Duration:} Full Year (36 weeks)\\
\textbf{Prerequisites:} Grade 3 or demonstrated readiness

\section{Course Description}

Grade 4 Social Studies typically focuses on state history, geography, economics, and government. Students explore their state's history, notable figures, geography, natural resources, and contribution to the United States.

\section{Course Goals}

Upon completion of this course, students will be able to:

\begin{itemize}
    \item Understand state history and important events
    \item Understand state geography and regions
    \item Understand state government and how it functions
    \item Understand the role of state government in students' lives
    \item Understand state culture and diversity
    \item Understand economic activities in the state
    \item Use maps and geographic tools
    \item Demonstrate citizenship in state and local contexts
\end{itemize}

\section{Assessment and Grading}

\begin{center}
\begin{tabular}{|l|c|}
\hline
\textbf{Assessment Category} & \textbf{Percentage} \\
\hline
Understanding of Concepts & 30\% \\
\hline
Projects and Presentations & 25\% \\
\hline
Map and Geography Skills & 20\% \\
\hline
Participation and Discussions & 15\% \\
\hline
Daily Work & 10\% \\
\hline
\end{tabular}
\end{center}

\section{Units of Study}

\subsection{Unit 1: State Geography and Regions (5-6 weeks)}

Topics:
\begin{itemize}
\item Understanding major landforms and water features of state/region
\item Reading and using various types of maps
\item Understanding map skills (scale, symbols, legends)
\item Exploring climate and natural resources
\item Understanding how geography affects settlement
\item Investigating regions of the state/country
\item Understanding latitude, longitude, and geographic features
\end{itemize}

\subsection{Unit 2: State or Local History (6-7 weeks)}

Topics:
\begin{itemize}
\item Understanding state/local historical events
\item Learning about significant historical figures
\item Exploring Native American history in the region
\item Understanding how communities developed
\item Investigating immigration and settlement patterns
\item Understanding cultural contributions
\item Creating timelines of important events
\end{itemize}

\subsection{Unit 3: Government and Citizenship (5-6 weeks)}

Topics:
\begin{itemize}
\item Understanding local, state, and national government structures
\item Learning about rights and responsibilities
\item Understanding rules, laws, and how they're made
\item Exploring how governments provide services
\item Understanding voting and democratic processes
\item Learning about civic participation
\item Understanding government leaders and their roles
\end{itemize}

\subsection{Unit 4: Economics and Trade (5-6 weeks)}

Topics:
\begin{itemize}
\item Understanding producers and consumers
\item Exploring specialization and trade
\item Understanding supply and demand
\item Investigating imports and exports
\item Understanding economic interdependence
\item Learning about natural resources and economic activities
\item Understanding scarcity and making choices
\end{itemize}

\subsection{Unit 5: Native American Cultures (4-5 weeks)}

Topics:
\begin{itemize}
\item Understanding Native American groups and regions
\item Exploring cultural traditions and practices
\item Understanding how environment shaped cultures
\item Learning about contributions and achievements
\item Understanding historical conflicts and impacts
\item Appreciating cultural diversity
\item Examining historical and contemporary Native American life
\end{itemize}

\subsection{Unit 6: Westward Expansion and Development (5-6 weeks)}

Topics:
\begin{itemize}
\item Understanding reasons for western expansion
\item Exploring the journey westward (if applicable to region)
\item Understanding impact on Native Americans
\item Learning about pioneers and settlements
\item Understanding technological developments
\item Investigating cultural changes and diversity
\item Understanding economic development
\end{itemize}

\subsection{Unit 7: American Symbols, Holidays, and Culture (4-5 weeks)}

Topics:
\begin{itemize}
\item Understanding American symbols and their meanings
\item Learning about national holidays and their significance
\item Understanding patriotism and national pride
\item Exploring cultural contributions and celebrations
\item Understanding American diversity and unity
\item Learning about important American documents
\item Understanding American values and ideals
\end{itemize}

\clearpage

\part{Grade 5}

\chapter{Grade 5 English/Language Arts}

\section{Course Overview}

\textbf{Grade Level:} Grade 5\\
\textbf{Duration:} Full Year (36 weeks)\\
\textbf{Prerequisites:} Grade 4 or demonstrated readiness

\section{Course Description}

Grade 5 English/Language Arts prepares students for more advanced literacy. Students read increasingly complex texts, analyze author's craft and structure, write for multiple purposes with sophistication, and apply conventions of standard English.

\section{Course Goals}

Upon completion of this course, students will be able to:

\begin{itemize}
    \item Quote accurately from texts to support analysis
    \item Determine theme and summarize key details
    \item Analyze character development and plot structure
    \item Compare multiple texts on similar topics
    \item Determine main idea and supporting details
    \item Write well-developed opinion, informative, and narrative texts
    \item Use varied sentence structures and precise language
    \item Use correct conventions and mechanics
\end{itemize}

\section{Assessment and Grading}

\begin{center}
\begin{tabular}{|l|c|}
\hline
\textbf{Assessment Category} & \textbf{Percentage} \\
\hline
Reading Analysis and Comprehension & 35\% \\
\hline
Writing & 40\% \\
\hline
Conventions and Language & 15\% \\
\hline
Participation & 10\% \\
\hline
\end{tabular}
\end{center}

\section{Units of Study}

\subsection{Unit 1: State Geography and Regions (5-6 weeks)}

Topics:
\begin{itemize}
\item Understanding major landforms and water features of state/region
\item Reading and using various types of maps
\item Understanding map skills (scale, symbols, legends)
\item Exploring climate and natural resources
\item Understanding how geography affects settlement
\item Investigating regions of the state/country
\item Understanding latitude, longitude, and geographic features
\end{itemize}

\subsection{Unit 2: State or Local History (6-7 weeks)}

Topics:
\begin{itemize}
\item Understanding state/local historical events
\item Learning about significant historical figures
\item Exploring Native American history in the region
\item Understanding how communities developed
\item Investigating immigration and settlement patterns
\item Understanding cultural contributions
\item Creating timelines of important events
\end{itemize}

\subsection{Unit 3: Government and Citizenship (5-6 weeks)}

Topics:
\begin{itemize}
\item Understanding local, state, and national government structures
\item Learning about rights and responsibilities
\item Understanding rules, laws, and how they're made
\item Exploring how governments provide services
\item Understanding voting and democratic processes
\item Learning about civic participation
\item Understanding government leaders and their roles
\end{itemize}

\subsection{Unit 4: Economics and Trade (5-6 weeks)}

Topics:
\begin{itemize}
\item Understanding producers and consumers
\item Exploring specialization and trade
\item Understanding supply and demand
\item Investigating imports and exports
\item Understanding economic interdependence
\item Learning about natural resources and economic activities
\item Understanding scarcity and making choices
\end{itemize}

\subsection{Unit 5: Native American Cultures (4-5 weeks)}

Topics:
\begin{itemize}
\item Understanding Native American groups and regions
\item Exploring cultural traditions and practices
\item Understanding how environment shaped cultures
\item Learning about contributions and achievements
\item Understanding historical conflicts and impacts
\item Appreciating cultural diversity
\item Examining historical and contemporary Native American life
\end{itemize}

\subsection{Unit 6: Westward Expansion and Development (5-6 weeks)}

Topics:
\begin{itemize}
\item Understanding reasons for western expansion
\item Exploring the journey westward (if applicable to region)
\item Understanding impact on Native Americans
\item Learning about pioneers and settlements
\item Understanding technological developments
\item Investigating cultural changes and diversity
\item Understanding economic development
\end{itemize}

\subsection{Unit 7: American Symbols, Holidays, and Culture (4-5 weeks)}

Topics:
\begin{itemize}
\item Understanding American symbols and their meanings
\item Learning about national holidays and their significance
\item Understanding patriotism and national pride
\item Exploring cultural contributions and celebrations
\item Understanding American diversity and unity
\item Learning about important American documents
\item Understanding American values and ideals
\end{itemize}

\chapter{Grade 5 Mathematics}

\section{Course Overview}

\textbf{Grade Level:} Grade 5\\
\textbf{Duration:} Full Year (36 weeks)\\
\textbf{Prerequisites:} Grade 4 or demonstrated readiness

\section{Course Description}

Grade 5 Mathematics develops fluency with fractions and decimals, extends multiplication and division to fractions, explores volume and geometric properties, and applies mathematical thinking to real-world problems.

\section{Course Goals}

Upon completion of this course, students will be able to:

\begin{itemize}
    \item Fluently add, subtract, multiply, and divide fractions
    \item Understand decimal place value and operations
    \item Multiply and divide whole numbers and decimals
    \item Solve problems involving fractions and decimals
    \item Understand volume and calculate it
    \item Classify shapes based on properties
    \item Represent and interpret data in various formats
    \item Apply mathematical thinking to real-world situations
\end{itemize}

\section{Assessment and Grading}

\begin{center}
\begin{tabular}{|l|c|}
\hline
\textbf{Assessment Category} & \textbf{Percentage} \\
\hline
Computational Fluency & 35\% \\
\hline
Problem-Solving & 30\% \\
\hline
Measurement and Geometry & 20\% \\
\hline
Participation & 10\% \\
\hline
Daily Work & 5\% \\
\hline
\end{tabular}
\end{center}

\section{Units of Study}

\subsection{Unit 1: Place Value and Whole Number Operations (4-5 weeks)}

Topics:
\begin{itemize}
\item Understanding place value to millions
\item Reading and writing large numbers
\item Comparing and ordering numbers
\item Rounding to various place values
\item Fluent addition and subtraction
\item Estimating sums and differences
\item Solving word problems with multiple steps
\end{itemize}

\subsection{Unit 2: Multiplication and Division (6-7 weeks)}

Topics:
\begin{itemize}
\item Multiplying multi-digit whole numbers
\item Dividing multi-digit whole numbers
\item Understanding remainders
\item Estimating products and quotients
\item Using standard algorithms
\item Solving multiplication and division word problems
\item Understanding factors and multiples
\end{itemize}

\subsection{Unit 3: Fractions and Decimal Connections (7-8 weeks)}

Topics:
\begin{itemize}
\item Understanding fractions as division
\item Identifying and naming fractions
\item Understanding equivalent fractions
\item Comparing and ordering fractions
\item Adding and subtracting fractions with like and unlike denominators
\item Multiplying fractions
\item Understanding decimals and fraction/decimal equivalence
\end{itemize}

\subsection{Unit 4: Decimal Operations (6-7 weeks)}

Topics:
\begin{itemize}
\item Reading and writing decimals
\item Comparing and ordering decimals
\item Rounding decimals
\item Adding and subtracting decimals
\item Multiplying and dividing decimals
\item Using decimals in money and measurement
\item Solving decimal word problems
\end{itemize}

\subsection{Unit 5: Measurement and Conversions (5-6 weeks)}

Topics:
\begin{itemize}
\item Measuring length, weight, volume, and time
\item Converting units within measurement systems
\item Understanding perimeter and area of rectangles
\item Understanding volume of rectangular prisms
\item Solving measurement word problems
\item Using appropriate measurement tools
\item Understanding metric and customary systems
\end{itemize}

\subsection{Unit 6: Geometry and Coordinate Planes (5-6 weeks)}

Topics:
\begin{itemize}
\item Classifying two-dimensional shapes
\item Understanding properties of polygons
\item Understanding angles and angle relationships
\item Graphing points on a coordinate plane
\item Understanding translations, reflections, and rotations
\item Understanding symmetry
\item Exploring congruence and similarity
\end{itemize}

\subsection{Unit 7: Data and Graphing (4-5 weeks)}

Topics:
\begin{itemize}
\item Collecting and organizing data
\item Creating and interpreting various graphs
\item Understanding measures of central tendency
\item Understanding probability basics
\item Creating and analyzing line plots
\item Drawing conclusions from data
\item Making predictions based on data
\end{itemize}

\subsection{Unit 8: Patterns and Algebraic Thinking (4-5 weeks)}

Topics:
\begin{itemize}
\item Identifying and extending patterns
\item Understanding number relationships
\item Understanding expressions and equations
\item Solving simple equations
\item Using variables
\item Understanding order of operations
\item Solving two-step word problems
\end{itemize}

\chapter{Grade 5 Science}

\section{Course Overview}

\textbf{Grade Level:} Grade 5\\
\textbf{Duration:} Full Year (36 weeks)\\
\textbf{Prerequisites:} Grade 4 or demonstrated readiness

\section{Course Description}

Grade 5 Science deepens understanding of life, physical, and earth systems. Students engage in guided inquiry, use scientific tools and technology, and develop sophisticated explanations of natural phenomena.

\section{Course Goals}

Upon completion of this course, students will be able to:

\begin{itemize}
    \item Conduct systematic scientific investigations
    \item Analyze and interpret data appropriately
    \item Develop explanations supported by evidence
    \item Understand matter and its properties
    \item Understand energy and transformations
    \item Understand life systems and interactions
    \item Understand earth systems and processes
    \item Apply scientific thinking to real-world problems
\end{itemize}

\section{Assessment and Grading}

\begin{center}
\begin{tabular}{|l|c|}
\hline
\textbf{Assessment Category} & \textbf{Percentage} \\
\hline
Investigations and Labs & 40\% \\
\hline
Data Analysis & 25\% \\
\hline
Scientific Reasoning & 20\% \\
\hline
Participation and Safety & 10\% \\
\hline
Science Journals & 5\% \\
\hline
\end{tabular}
\end{center}

\section{Units of Study}

\subsection{Unit 1: Human Body Systems (7-8 weeks)}

Topics:
\begin{itemize}
\item Understanding body systems (skeletal, muscular, digestive, circulatory, respiratory)
\item Investigating organs and their functions
\item Understanding how systems work together
\item Learning about nutrition and healthy eating
\item Exploring exercise and fitness
\item Understanding disease and immunity
\item Learning about growth and development
\end{itemize}

Investigations:
\begin{itemize}
\item Modeling body systems
\item Investigating pulse and heart rate
\item Analyzing nutrition labels
\item Studying effects of exercise
\end{itemize}

\subsection{Unit 2: Cell Structure and Life Processes (6-7 weeks)}

Topics:
\begin{itemize}
\item Understanding cell structure and organelles
\item Comparing plant and animal cells
\item Understanding photosynthesis and respiration
\item Investigating cell division
\item Understanding growth and reproduction
\item Exploring how cells function
\item Understanding living and non-living things
\end{itemize}

Investigations:
\begin{itemize}
\item Observing cells with microscopes
\item Modeling cell structures
\item Investigating photosynthesis
\end{itemize}

\subsection{Unit 3: Heredity and Genetics (6-7 weeks)}

Topics:
\begin{itemize}
\item Understanding DNA and genes
\item Exploring heredity and traits
\item Understanding dominant and recessive traits
\item Investigating inheritance patterns
\item Understanding variation within species
\item Exploring mutations
\item Understanding genetic engineering basics
\end{itemize}

Investigations:
\begin{itemize}
\item Studying trait inheritance in families
\item Modeling DNA
\item Exploring Punnett squares
\end{itemize}

\subsection{Unit 4: Evolution and Natural Selection (6-7 weeks)}

Topics:
\begin{itemize}
\item Understanding evidence of evolution
\item Exploring fossil record
\item Understanding natural selection
\item Investigating adaptations and survival
\item Understanding speciation
\item Exploring classification of organisms
\item Understanding extinction
\end{itemize}

Investigations:
\begin{itemize}
\item Analyzing fossil records
\item Studying adaptations in local organisms
\item Investigating camouflage and adaptation
\end{itemize}

\subsection{Unit 5: Ecosystems and Biomes (6-7 weeks)}

Topics:
\begin{itemize}
\item Understanding biotic and abiotic factors
\item Exploring energy flow through food webs
\item Understanding cycling of matter
\item Investigating different biomes
\item Understanding population dynamics
\item Exploring competition and predation
\item Understanding succession
\end{itemize}

Investigations:
\begin{itemize}
\item Creating ecosystem models
\item Studying local food webs
\item Investigating population interactions
\item Observing ecosystem changes
\end{itemize}

\subsection{Unit 6: Space and Earth Systems (7-8 weeks)}

Topics:
\begin{itemize}
\item Understanding the solar system
\item Exploring planets and their characteristics
\item Understanding Earth's movements and seasons
\item Investigating the moon's phases
\item Understanding gravity and orbits
\item Exploring the universe
\item Understanding telescopes and space exploration
\end{itemize}

Investigations:
\begin{itemize}
\item Modeling planetary motion
\item Tracking moon phases
\item Using star maps
\item Building telescopes
\end{itemize}

\subsection{Unit 7: Earth's Structure and Processes (6-7 weeks)}

Topics:
\begin{itemize}
\item Understanding Earth's layers
\item Investigating plate tectonics
\item Understanding earthquakes and volcanoes
\item Exploring mineral and rock formation
\item Understanding weathering, erosion, and deposition
\item Investigating fossil formation
\item Understanding geological time
\end{itemize}

Investigations:
\begin{itemize}
\item Modeling Earth's layers
\item Investigating plate movements
\item Creating mineral and rock models
\item Studying geologic time scales
\end{itemize}

\chapter{Grade 5 Social Studies}

\section{Course Overview}

\textbf{Grade Level:} Grade 5\\
\textbf{Duration:} Full Year (36 weeks)\\
\textbf{Prerequisites:} Grade 4 or demonstrated readiness

\section{Course Description}

Grade 5 Social Studies typically focuses on United States history and geography. Students explore American history from pre-Columbian times through the present, understand geographic features and regions, and develop civic knowledge.

\section{Course Goals}

Upon completion of this course, students will be able to:

\begin{itemize}
    \item Understand major eras and events in American history
    \item Understand the geographic features of the United States
    \item Understand the development of American government
    \item Understand the causes and effects of historical events
    \item Compare perspectives of different historical figures
    \item Understand cultural contributions and diversity
    \item Use maps and analyze geographic information
    \item Demonstrate civic understanding and values
\end{itemize}

\section{Assessment and Grading}

\begin{center}
\begin{tabular}{|l|c|}
\hline
\textbf{Assessment Category} & \textbf{Percentage} \\
\hline
Historical Understanding & 30\% \\
\hline
Projects and Presentations & 25\% \\
\hline
Analysis and Critical Thinking & 20\% \\
\hline
Map and Geography Skills & 15\% \\
\hline
Participation & 10\% \\
\hline
\end{tabular}
\end{center}

\section{Units of Study}

\subsection{Unit 1: Early Civilizations (6-7 weeks)}

Topics:
\begin{itemize}
\item Understanding early human societies
\item Exploring development of agriculture
\item Investigating early civilizations (Mesopotamia, Egypt, Indus Valley, China)
\item Understanding government and religion in ancient societies
\item Exploring trade and cultural exchange
\item Understanding technological advances
\item Learning about social structures
\end{itemize}

\subsection{Unit 2: Classical Civilizations (6-7 weeks)}

Topics:
\begin{itemize}
\item Exploring classical Greece and Rome
\item Understanding democracy and government systems
\item Investigating cultural achievements
\item Learning about military and conquest
\item Exploring trade networks
\item Understanding social hierarchies
\item Investigating lasting influences
\end{itemize}

\subsection{Unit 3: Major World Religions (5-6 weeks)}

Topics:
\begin{itemize}
\item Understanding major religions (Christianity, Islam, Judaism, Hinduism, Buddhism)
\item Exploring origins and key figures
\item Understanding beliefs and practices
\item Investigating historical influences
\item Understanding religious contributions to society
\item Exploring diversity within religions
\item Understanding religious tolerance
\end{itemize}

\subsection{Unit 4: Medieval Period (5-6 weeks)}

Topics:
\begin{itemize}
\item Understanding feudalism and social structure
\item Exploring kingdoms and empires
\item Investigating cultural and technological developments
\item Understanding the Crusades
\item Exploring trade and exploration
\item Understanding transitions to modern period
\end{itemize}

\subsection{Unit 5: Age of Exploration and Colonization (6-7 weeks)}

Topics:
\begin{itemize}
\item Understanding causes of exploration
\item Investigating early explorers and routes
\item Understanding colonization and settlement
\item Exploring cultural contact and exchange
\item Understanding impact on indigenous peoples
\item Investigating economic systems (mercantilism, slavery)
\item Understanding lasting consequences
\end{itemize}

\subsection{Unit 6: Development of Democratic Ideals (5-6 weeks)}

Topics:
\begin{itemize}
\item Understanding Enlightenment ideas
\item Investigating American Revolution
\item Understanding founding documents
\item Exploring rights and freedoms
\item Understanding government structure
\item Investigating other revolutions
\item Understanding spread of democratic ideals
\end{itemize}

\subsection{Unit 7: Industrial Revolution and Global Impact (5-6 weeks)}

Topics:
\begin{itemize}
\item Understanding causes and effects of industrialization
\item Exploring technological innovations
\item Investigating economic changes
\item Understanding social impacts
\item Exploring global expansion
\item Understanding lasting changes
\item Investigating rise of modern nations
\end{itemize}

\clearpage

\part{Grade 6 (Middle School)}

\chapter{Grade 6 English/Language Arts}

\section{Course Overview}

\textbf{Grade Level:} Grade 6\\
\textbf{Duration:} Full Year (36 weeks)\\
\textbf{Prerequisites:} Grade 5 or demonstrated readiness

\section{Course Description}

Grade 6 English/Language Arts prepares students for high school-level literacy. Students read and analyze complex texts, develop sophisticated writing skills across multiple genres, and deepen understanding of literary elements and author's craft.

\section{Course Goals}

Upon completion of this course, students will be able to:

\begin{itemize}
    \item Analyze texts for theme, character development, and author's craft
    \item Determine central ideas and summarize complex texts
    \item Analyze multiple interpretations of texts
    \item Write well-developed essays and creative pieces
    \item Use varied sentence structures and precise vocabulary
    \item Analyze the author's purpose and point of view
    \item Evaluate arguments and evidence
    \item Demonstrate command of grammar and conventions
\end{itemize}

\section{Assessment and Grading}

\begin{center}
\begin{tabular}{|l|c|}
\hline
\textbf{Assessment Category} & \textbf{Percentage} \\
\hline
Literary Analysis and Reading & 35\% \\
\hline
Writing and Composition & 40\% \\
\hline
Grammar and Language & 15\% \\
\hline
Participation & 10\% \\
\hline
\end{tabular}
\end{center}

\section{Units of Study}

\subsection{Unit 1: Literary Analysis and Interpretation (10-12 weeks)}

Topics:
\begin{itemize}
\item Analyzing complex characters and character development
\item Understanding universal themes and messages
\item Analyzing plot structure including complex conflicts
\item Understanding point of view and narrative perspective
\item Analyzing figurative language (metaphor, simile, personification, idiom)
\item Making inferences and supporting with textual evidence
\item Understanding author's purpose and craft
\item Comparing texts across genres and cultures
\end{itemize}

\subsection{Unit 2: Informational Reading and Analysis (8-10 weeks)}

Topics:
\begin{itemize}
\item Identifying central ideas and supporting details
\item Analyzing text structure and organization
\item Understanding author's perspective and bias
\item Evaluating credibility and reliability of sources
\item Analyzing relationships between ideas
\item Drawing conclusions and synthesizing information
\item Using graphics and multimedia
\item Understanding media techniques and persuasion
\end{itemize}

\subsection{Unit 3: Writing for Different Purposes (Ongoing)}

Duration: Integrated throughout the year

Topics:
\begin{itemize}
\item Writing arguments with claims and evidence
\item Writing informative/explanatory texts with research
\item Writing narratives with dialogue and descriptive details
\item Using organizational strategies
\item Developing and supporting thesis statements
\item Citing sources and avoiding plagiarism
\item Revising for clarity, organization, style, and voice
\item Publishing for different audiences
\end{itemize}

\subsection{Unit 4: Communication and Presentation (Ongoing)}

Duration: Integrated throughout the year

Topics:
\begin{itemize}
\item Engaging in substantive discussions
\item Evaluating speakers' arguments and reasoning
\item Presenting information persuasively
\item Using multimedia and visual aids effectively
\item Practicing active listening
\item Using standard English in formal settings
\item Giving and receiving constructive feedback
\end{itemize}

\subsection{Unit 5: Language Development and Conventions (Ongoing)}

Duration: Integrated throughout the year

Topics:
\begin{itemize}
\item Understanding complex sentence structures
\item Using pronouns correctly
\item Understanding and using verb tenses and moods
\item Using correct punctuation and capitalization
\item Understanding verbals and phrases
\item Using coordination and subordination
\item Maintaining consistency
\item Recognizing and correcting errors
\end{itemize}

\subsection{Unit 6: Vocabulary and Academic Discourse (Ongoing)}

Duration: Integrated throughout the year

Topics:
\begin{itemize}
\item Using context clues and academic vocabulary
\item Understanding word families and etymology
\item Using morphological analysis
\item Building discipline-specific vocabulary
\item Understanding connotation, denotation, and nuance
\item Recognizing relationships between words
\item Using reference materials strategically
\item Building vocabulary across disciplines
\end{itemize}


\chapter{Grade 6 Mathematics}

\section{Course Overview}

\textbf{Grade Level:} Grade 6\\
\textbf{Duration:} Full Year (36 weeks)\\
\textbf{Prerequisites:} Grade 5 or demonstrated readiness

\section{Course Description}

Grade 6 Mathematics extends understanding of rational numbers, develops algebraic thinking, explores ratios and proportional relationships, and applies mathematical concepts to real-world situations.

\section{Course Goals}

Upon completion of this course, students will be able to:

\begin{itemize}
    \item Fluently divide fractions by fractions
    \item Compute with multi-digit decimals
    Understand positive and negative integers
    \item Understand and use variables in expressions
    \item Understand ratios and use them to solve problems
    \item Understand percent and use it in applications
    \item Analyze and display data
    \item Understand geometric properties and measurements
\end{itemize}

\section{Assessment and Grading}

\begin{center}
\begin{tabular}{|l|c|}
\hline
\textbf{Assessment Category} & \textbf{Percentage} \\
\hline
Computational Skills & 30\% \\
\hline
Problem-Solving and Reasoning & 35\% \\
\hline
Data and Geometry & 20\% \\
\hline
Participation & 10\% \\
\hline
Daily Work & 5\% \\
\hline
\end{tabular}
\end{center}

\section{Units of Study}

\subsection{Unit 1: Whole Numbers and Rational Numbers (5-6 weeks)}

Topics:
\begin{itemize}
\item Understanding factors, multiples, and primes
\item Using order of operations
\item Understanding fractions and decimals
\item Converting between fractions and decimals
\item Comparing and ordering rational numbers
\item Understanding absolute value
\item Identifying patterns in sequences
\end{itemize}

\subsection{Unit 2: Operations with Fractions and Decimals (6-7 weeks)}

Topics:
\begin{itemize}
\item Adding and subtracting fractions
\item Multiplying and dividing fractions
\item Adding, subtracting, multiplying, and dividing decimals
\item Estimating results
\item Solving word problems with fractions and decimals
\item Understanding relationships between operations
\item Solving multi-step problems
\end{itemize}

\subsection{Unit 3: Ratios, Rates, and Proportions (6-7 weeks)}

Topics:
\begin{itemize}
\item Understanding ratios and rates
\item Comparing unit rates
\item Understanding proportional relationships
\item Solving proportion problems
\item Understanding percent
\item Calculating discounts, tax, and interest
\item Solving real-world problems with ratios
\end{itemize}

\subsection{Unit 4: Expressions and Equations (6-7 weeks)}

Topics:
\begin{itemize}
\item Writing and evaluating expressions
\item Understanding variables
\item Combining like terms
\item Writing and solving one-step equations
\item Writing and solving two-step equations
\item Understanding solving strategies
\item Solving word problems with equations
\end{itemize}

\subsection{Unit 5: Inequalities and Graphing (4-5 weeks)}

Topics:
\begin{itemize}
\item Writing and graphing inequalities
\item Understanding solution sets
\item Solving one-step inequalities
\item Graphing on coordinate planes
\item Understanding independent and dependent variables
\item Analyzing relationships in graphs
\item Making predictions from graphs
\end{itemize}

\subsection{Unit 6: Geometry (5-6 weeks)}

Topics:
\begin{itemize}
\item Understanding angles and angle relationships
\item Understanding triangles and quadrilaterals
\item Finding area of triangles and quadrilaterals
\item Finding perimeter and area of composite figures
\item Understanding volume and surface area
\item Understanding congruence and similarity
\item Understanding transformations
\end{itemize}

\subsection{Unit 7: Statistics and Probability (5-6 weeks)}

Topics:
\begin{itemize}
\item Collecting and organizing data
\item Creating and interpreting graphs
\item Understanding mean, median, mode, and range
\item Understanding variability
\item Understanding probability
\item Calculating theoretical and experimental probability
\item Making predictions based on probability
\end{itemize}

\chapter{Grade 6 Science}

\section{Course Overview}

\textbf{Grade Level:} Grade 6\\
\textbf{Duration:} Full Year (36 weeks)\\
\textbf{Prerequisites:} Grade 5 or demonstrated readiness

\section{Course Description}

Grade 6 Science explores earth and life systems in greater depth. Students conduct inquiry-based investigations, use scientific tools and technology, and develop evidence-based explanations of natural phenomena at the middle school level.

\section{Course Goals}

Upon completion of this course, students will be able to:

\begin{itemize}
    \item Design and conduct systematic investigations
    \item Collect and analyze data using appropriate methods
    \item Develop explanations based on evidence
    \item Understand earth's structure and processes
    \item Understand the water cycle and weather
    \item Understand plate tectonics
    \item Understand organisms and populations
    \item Understand ecosystems and their interactions
\end{itemize}

\section{Assessment and Grading}

\begin{center}
\begin{tabular}{|l|c|}
\hline
\textbf{Assessment Category} & \textbf{Percentage} \\
\hline
Investigations and Labs & 40\% \\
\hline
Data Analysis and Conclusions & 25\% \\
\hline
Scientific Reasoning & 20\% \\
\hline
Participation and Safety & 10\% \\
\hline
Science Journals & 5\% \\
\hline
\end{tabular}
\end{center}

\section{Units of Study}

\subsection{Unit 1: Structure and Function of Life (7-8 weeks)}

Topics:
\begin{itemize}
\item Understanding organization of life (cells, tissues, organs, systems)
\item Investigating plant and animal structures
\item Understanding how structures relate to function
\item Exploring photosynthesis and cellular respiration
\item Understanding body systems and their interactions
\item Investigating reproduction and growth
\item Understanding classification of organisms
\end{itemize}

\subsection{Unit 2: Heredity and Evolution (7-8 weeks)}

Topics:
\begin{itemize}
\item Understanding inheritance and traits
\item Exploring DNA and genetics
\item Understanding natural selection
\item Investigating adaptation and evolution
\item Understanding evidence of evolution
\item Exploring fossil records
\item Understanding diversity of life
\end{itemize}

\subsection{Unit 3: Ecosystems and Biotic Interactions (6-7 weeks)}

Topics:
\begin{itemize}
\item Understanding biotic and abiotic factors
\item Exploring food webs and energy flow
\item Understanding nutrient cycles
\item Investigating population interactions
\item Understanding succession
\item Exploring biome characteristics
\item Understanding human impacts on ecosystems
\end{itemize}

\subsection{Unit 4: Forces and Motion (6-7 weeks)}

Topics:
\begin{itemize}
\item Understanding Newton's laws of motion
\item Investigating friction and resistance
\item Understanding velocity and acceleration
\item Exploring forces and balanced/unbalanced forces
\item Understanding gravitational force
\item Investigating simple machines
\item Exploring work and energy
\end{itemize}

\subsection{Unit 5: Energy Transfer and Transformation (6-7 weeks)}

Topics:
\begin{itemize}
\item Understanding forms of energy
\item Investigating energy transfer
\item Understanding conservation of energy
\item Exploring heat and temperature
\item Understanding thermal energy transfer
\item Investigating light and sound
\item Understanding electromagnetic waves
\end{itemize}

\subsection{Unit 6: Matter and Chemical Reactions (6-7 weeks)}

Topics:
\begin{itemize}
\item Understanding atomic structure
\item Exploring periodic table
\item Understanding chemical bonding
\item Investigating physical and chemical properties
\item Understanding chemical reactions
\item Exploring states of matter
\item Understanding conservation of mass
\end{itemize}

\subsection{Unit 7: Earth Systems (6-7 weeks)}

Topics:
\begin{itemize}
\item Understanding Earth's structure and composition
\item Investigating plate tectonics
\item Exploring weather and climate
\item Understanding water cycle and weather patterns
\item Investigating rocks and minerals
\item Understanding weathering and erosion
\item Exploring natural resources and conservation
\end{itemize}

\chapter{Grade 6 Social Studies}

\section{Course Overview}

\textbf{Grade Level:} Grade 6\\
\textbf{Duration:} Full Year (36 weeks)\\
\textbf{Prerequisites:} Grade 5 or demonstrated readiness

\section{Course Description}

Grade 6 Social Studies focuses on world history, geography, and cultures. Students explore ancient civilizations, world geography, cultural diversity, and the development of major world religions and governments.

\section{Course Goals}

Upon completion of this course, students will be able to:

\begin{itemize}
    \item Understand major ancient civilizations and their contributions
    \item Analyze causes and effects of historical events
    \item Understand world geography and its influence on civilizations
    \item Understand world cultures and religions
    \item Analyze different governmental systems
    \item Understand economic systems and trade
    \item Compare and contrast historical perspectives
    \item Apply historical thinking to current issues
\end{itemize}

\section{Assessment and Grading}

\begin{center}
\begin{tabular}{|l|c|}
\hline
\textbf{Assessment Category} & \textbf{Percentage} \\
\hline
Historical Understanding & 30\% \\
\hline
Analysis and Critical Thinking & 25\% \\
\hline
Projects and Presentations & 20\% \\
\hline
Geography and Map Skills & 15\% \\
\hline
Participation & 10\% \\
\hline
\end{tabular}
\end{center}

\section{Units of Study}

\subsection{Unit 1: Ancient Civilizations (6-7 weeks)}

Topics:
\begin{itemize}
\item Understanding human prehistory and development
\item Exploring Mesopotamian civilization
\item Investigating Egyptian civilization
\item Understanding Indus Valley and early China
\item Exploring contributions of ancient civilizations
\item Understanding development of writing and government
\item Investigating trade and cultural exchange
\end{itemize}

\subsection{Unit 2: Classical Civilizations and Empires (6-7 weeks)}

Topics:
\begin{itemize}
\item Exploring classical Greece
\item Understanding Roman Republic and Empire
\item Investigating Hellenistic period
\item Understanding government and political systems
\item Exploring cultural achievements and philosophy
\item Understanding military conquests and empires
\item Investigating influences on Western civilization
\end{itemize}

\subsection{Unit 3: World Religions (6-7 weeks)}

Topics:
\begin{itemize}
\item Understanding origins and core beliefs of major world religions
\item Exploring Buddhism, Christianity, Hinduism, Islam, Judaism
\item Understanding religious practices and traditions
\item Investigating historical influences of religions
\item Understanding religious architecture and art
\item Exploring religious texts
\item Understanding religious diversity and tolerance
\end{itemize}

\subsection{Unit 4: Medieval Europe and Beyond (6-7 weeks)}

Topics:
\begin{itemize}
\item Understanding feudalism and manorialism
\item Exploring medieval kingdoms and government
\item Investigating Byzantine Empire
\item Understanding Islamic Golden Age
\item Exploring medieval culture and achievements
\item Understanding crusades and cultural contacts
\item Investigating medieval Africa and Asia
\end{itemize}

\subsection{Unit 5: Renaissance and Exploration (6-7 weeks)}

Topics:
\begin{itemize}
\item Understanding causes and characteristics of Renaissance
\item Exploring Renaissance art and architecture
\item Understanding humanist movement
\item Investigating Age of Exploration
\item Understanding colonization and colonialism
\item Exploring columbian exchange
\item Understanding global connections
\end{itemize}

\subsection{Unit 6: Early Modern Period (5-6 weeks)}

Topics:
\begin{itemize}
\item Understanding Scientific Revolution
\item Exploring Enlightenment ideas
\item Investigating religious reformations
\item Understanding absolute monarchies
\item Exploring mercantilism and early capitalism
\item Understanding technology and innovation
\item Investigating global trade networks
\end{itemize}

\clearpage

\part{Grades 7-8}

\chapter{Grade 7 English/Language Arts}

\section{Course Overview}

\textbf{Grade Level:} Grade 7\\
\textbf{Duration:} Full Year (36 weeks)\\
\textbf{Prerequisites:} Grade 6 or demonstrated readiness

\section{Course Description}

Grade 7 English/Language Arts develops sophisticated readers and writers. Students analyze complex texts, understand literary elements in depth, write persuasive and analytical essays, and demonstrate command of academic language.

\section{Course Goals}

Upon completion of this course, students will be able to:

\begin{itemize}
    \item Analyze literature for themes, character motivation, and symbolism
    \item Determine central ideas in complex texts
    \item Evaluate the credibility and relevance of evidence
    \item Write persuasive essays with logical arguments
    \item Write analytical essays examining text elements
    \item Use varied sentence structures for effect
    \item Demonstrate command of grammar and mechanics
    \item Engage in academic discussions and debates
\end{itemize}

\section{Required Materials}

\begin{itemize}
    \item Literature anthology
    \item Young adult novels
    \item Research and reference materials
    \item Writing guides and rubrics
\end{itemize}

\section{Assessment and Grading}

\begin{center}
\begin{tabular}{|l|c|}
\hline
\textbf{Assessment Category} & \textbf{Percentage} \\
\hline
Literary Analysis & 35\% \\
\hline
Writing Assignments & 40\% \\
\hline
Conventions and Language & 15\% \\
\hline
Participation and Discussion & 10\% \\
\hline
\end{tabular}
\end{center}

\section{Units of Study}

\subsection{Unit 1: Literature Analysis and Close Reading (10-12 weeks)}

Topics:
\begin{itemize}
\item Analyzing complex character motivation and development
\item Understanding universal themes and cultural context
\item Analyzing plot structure including exposition, rising action, climax, falling action
\item Understanding point of view, narrator reliability, and perspective
\item Analyzing figurative language and symbolism
\item Making inferences with textual support
\item Understanding author's purpose, craft, and style
\item Comparing thematic elements across texts
\item Analyzing tone and mood in literature
\end{itemize}

\subsection{Unit 2: Non-fiction Reading and Analysis (8-10 weeks)}

Topics:
\begin{itemize}
\item Identifying thesis, main ideas, and supporting evidence
\item Analyzing text structure (chronological, cause-effect, problem-solution, comparison)
\item Understanding author's perspective, bias, and purpose
\item Evaluating credibility, reliability, and validity
\item Distinguishing fact from opinion and inference
\item Analyzing relationships and connections between ideas
\item Synthesizing information from multiple sources
\item Understanding rhetorical devices and persuasive techniques
\item Evaluating media messages and advertising
\end{itemize}

\subsection{Unit 3: Argumentative Writing (Ongoing)}

Duration: Integrated throughout the year

Topics:
\begin{itemize}
\item Developing clear claims and strong thesis statements
\item Supporting claims with relevant evidence and reasoning
\item Addressing opposing viewpoints and counterarguments
\item Using logical reasoning and valid arguments
\item Organizing arguments effectively
\item Using credible sources and citations
\item Revising for logic, clarity, and persuasiveness
\item Publishing arguments for various audiences
\item Understanding fallacies and logical errors
\end{itemize}

\subsection{Unit 4: Informative Writing and Research (Ongoing)}

Duration: Integrated throughout the year

Topics:
\begin{itemize}
\item Conducting research using multiple sources
\item Organizing information effectively
\item Writing clear topic sentences and transitions
\item Integrating evidence, facts, and details
\item Properly citing sources
\item Avoiding plagiarism
\item Summarizing and paraphrasing
\item Creating works cited pages
\item Revising for accuracy and clarity
\end{itemize}

\subsection{Unit 5: Narrative Writing and Personal Expression (Ongoing)}

Duration: Integrated throughout the year

Topics:
\begin{itemize}
\item Writing engaging narratives with clear plot structure
\item Developing characters with dialogue and description
\item Using sensory details and vivid language
\item Varying sentence structure and pacing
\item Using reflection and insight
\item Understanding perspective and voice
\item Revising for effectiveness and impact
\item Experimenting with narrative techniques
\end{itemize}

\subsection{Unit 6: Speaking, Listening, and Discussion (Ongoing)}

Duration: Integrated throughout the year

Topics:
\begin{itemize}
\item Engaging in academic discussions
\item Posing and responding to questions
\item Evaluating speakers' arguments and reasoning
\item Making presentations and speeches
\item Using multimedia and visual aids
\item Practicing active listening
\item Providing constructive feedback
\item Using formal English in appropriate contexts
\end{itemize}

\chapter{Grade 7 Mathematics}

\section{Course Overview}

\textbf{Grade Level:} Grade 7\\
\textbf{Duration:} Full Year (36 weeks)\\
\textbf{Prerequisites:} Grade 6 or demonstrated readiness

\section{Course Description}

Grade 7 Mathematics develops algebraic thinking, explores ratios and proportional relationships in depth, works with rational numbers, develops understanding of geometry and measurement, and applies mathematics to real-world contexts.

\section{Course Goals}

Upon completion of this course, students will be able to:

\begin{itemize}
    \item Develop fluency with rational number operations
    \item Analyze proportional relationships
    \item Use rates and ratios in problem-solving
    \item Solve equations and inequalities
    \item Apply geometry concepts to real-world problems
    \item Analyze and display statistical data
    \item Calculate probability
    \item Apply mathematical thinking to complex problems
\end{itemize}

\section{Units of Study}

\textbf{Major topics include:}
\begin{itemize}
    \item Rational Numbers and Operations
    \item Expressions and Equations
    \item Ratios and Proportional Relationships
    \item Geometry and Measurement
    \item Statistics and Probability
\end{itemize}

\section{Assessment and Grading}

\begin{center}
\begin{tabular}{|l|c|}
\hline
\textbf{Assessment Category} & \textbf{Percentage} \\
\hline
Computational Skills & 30\% \\
\hline
Problem-Solving and Reasoning & 35\% \\
\hline
Geometry and Statistics & 20\% \\
\hline
Participation & 10\% \\
\hline
Daily Work & 5\% \\
\hline
\end{tabular}
\end{center}

\section{Units of Study}

\subsection{Unit 1: Ratios, Rates, and Proportional Relationships (6–7 weeks)}

Topics:
\begin{itemize}
\item Understanding ratios and equivalent ratios
\item Computing unit rates associated with ratios
\item Representing proportional relationships in tables and graphs
\item Identifying the constant of proportionality
\item Writing equations to represent proportional relationships
\item Solving real-world percent problems (tax, discount, markup, tip, percent error)
\end{itemize}

\subsection{Unit 2: Operations with Rational Numbers (6–7 weeks)}

Topics:
\begin{itemize}
\item Adding and subtracting integers on the number line
\item Multiplying and dividing integers
\item Adding, subtracting, multiplying, and dividing rational numbers
\item Understanding opposites, absolute value, and distance
\item Solving real-world problems with rational numbers
\item Using properties of operations to justify steps in computation
\end{itemize}

\subsection{Unit 3: Expressions, Equations, and Inequalities (6–7 weeks)}

Topics:
\begin{itemize}
\item Writing and evaluating algebraic expressions
\item Using properties to generate equivalent expressions
\item Solving multi-step equations with rational coefficients
\item Solving and graphing inequalities on a number line
\item Modeling word problems with equations and inequalities
\item Interpreting solutions in context
\end{itemize}

\subsection{Unit 4: Geometry – Scale Drawings and Angle Relationships (5–6 weeks)}

Topics:
\begin{itemize}
\item Interpreting and creating scale drawings
\item Using scale to compute actual lengths and areas
\item Identifying complementary, supplementary, vertical, and adjacent angles
\item Using angle relationships to find unknown angle measures
\item Solving real-world problems with scale drawings and angles
\end{itemize}

\subsection{Unit 5: Geometry – Area, Surface Area, and Volume (5–6 weeks)}

Topics:
\begin{itemize}
\item Finding area of triangles, quadrilaterals, and polygons
\item Finding surface area of prisms and pyramids
\item Finding volume of prisms and cylinders
\item Solving real-world problems involving area, surface area, and volume
\item Relating nets to surface area
\end{itemize}

\subsection{Unit 6: Statistics – Sampling and Inference (5–6 weeks)}

Topics:
\begin{itemize}
\item Understanding random sampling
\item Using samples to estimate population characteristics
\item Understanding variability in samples
\item Comparing two populations using data
\item Constructing and interpreting dot plots, histograms, and box plots
\item Informally assessing overlap of distributions
\end{itemize}

\subsection{Unit 7: Probability (4–5 weeks)}

Topics:
\begin{itemize}
\item Understanding probability as a number between 0 and 1
\item Finding experimental and theoretical probability
\item Using probability to predict relative frequency
\item Representing sample spaces (lists, tables, tree diagrams)
\item Understanding compound events
\item Designing and using simulations to estimate probabilities
\end{itemize}


\chapter{Grade 7 Science}

\section{Course Overview}

\textbf{Grade Level:} Grade 7\\
\textbf{Duration:} Full Year (36 weeks)\\
\textbf{Prerequisites:} Grade 6 or demonstrated readiness

\section{Course Description}

Grade 7 Science integrates life, physical, and earth science. Students conduct inquiry-based investigations, develop scientific explanations, and apply scientific thinking to real-world issues.

\section{Course Goals}

Upon completion of this course, students will be able to:

\begin{itemize}
    \item Design and conduct fair scientific investigations
    \item Analyze data and develop conclusions
    \item Construct evidence-based explanations
    \item Understand matter and energy transformations
    \item Understand forces and motion
    \item Understand waves and energy transfer
    \item Understand organisms and inheritance
    \item Understand ecosystems and evolution
    \item Understand earth systems and resources
\end{itemize}

\section{Units of Study}

\textbf{Major topics include:}
\begin{itemize}
    \item Matter and Properties
    \item Forces and Motion
    \item Energy and Waves
    \item Heredity and Inheritance
    \item Ecosystems and Evolution
    \item Earth Systems
\end{itemize}

\section{Assessment and Grading}

\begin{center}
\begin{tabular}{|l|c|}
\hline
\textbf{Assessment Category} & \textbf{Percentage} \\
\hline
Investigations and Labs & 40\% \\
\hline
Data Analysis & 25\% \\
\hline
Scientific Reasoning & 20\% \\
\hline
Participation and Safety & 10\% \\
\hline
Science Journals & 5\% \\
\hline
\end{tabular}
\end{center}

\section{Units of Study}

\subsection{Unit 1: Structure and Function in Living Systems (6–7 weeks)}

Topics:
\begin{itemize}
\item Cell theory and characteristics of living things
\item Structure and function of plant and animal cells
\item Levels of organization (cells, tissues, organs, systems)
\item Major human body systems and their interactions
\item Structure–function relationships in organisms
\end{itemize}

\subsection{Unit 2: Genetics and Heredity (6–7 weeks)}

Topics:
\begin{itemize}
\item Traits, genes, and alleles
\item Dominant and recessive inheritance patterns
\item Punnett squares and probability of traits
\item Asexual vs. sexual reproduction
\item Variations within a species
\item Ethical issues related to genetics (introductory)
\end{itemize}

\subsection{Unit 3: Evolution and Diversity of Life (5–6 weeks)}

Topics:
\begin{itemize}
\item Evidence for evolution (fossils, anatomy, DNA, embryos)
\item Natural selection and adaptation
\item Survival and reproduction in changing environments
\item Speciation and extinction
\item Organization and classification of organisms
\end{itemize}

\subsection{Unit 4: Ecosystems, Matter, and Energy (6–7 weeks)}

Topics:
\begin{itemize}
\item Biotic and abiotic factors in ecosystems
\item Food chains, food webs, and energy pyramids
\item Cycling of matter (water, carbon, nitrogen)
\item Population dynamics and carrying capacity
\item Human impacts on ecosystems and biodiversity
\end{itemize}

\subsection{Unit 5: Forces and Motion (6–7 weeks)}

Topics:
\begin{itemize}
\item Describing motion using distance, time, and speed
\item Balanced and unbalanced forces
\item Newton’s laws of motion (qualitative)
\item Friction, gravity, and normal force
\item Interpreting motion graphs
\end{itemize}

\subsection{Unit 6: Energy and Energy Transfer (5–6 weeks)}

Topics:
\begin{itemize}
\item Forms of energy (kinetic, potential, thermal, chemical, electrical)
\item Conservation of energy in systems
\item Energy transformations
\item Heat transfer (conduction, convection, radiation)
\item Simple machines and mechanical advantage (review/extension)
\end{itemize}

\subsection{Unit 7: Waves, Sound, and Light (5–6 weeks)}

Topics:
\begin{itemize}
\item Characteristics of waves (wavelength, frequency, amplitude)
\item Mechanical vs. electromagnetic waves
\item Properties of sound waves
\item Properties of light; reflection and refraction
\item Applications of waves in technology (basic)
\end{itemize}


\chapter{Grade 7 Social Studies}

\section{Course Overview}

\textbf{Grade Level:} Grade 7\\
\textbf{Duration:} Full Year (36 weeks)\\
\textbf{Prerequisites:} Grade 6 or demonstrated readiness

\section{Course Description}

Grade 7 Social Studies focuses on world history and geography, exploring the major world regions, cultures, and historical developments.

\section{Course Goals}

Upon completion of this course, students will be able to:

\begin{itemize}
    \item Understand major periods and events in world history
    \item Understand world geography and its influence on societies
    \item Analyze causes and effects of historical change
    \item Understand diverse cultures and religions
    \item Compare governmental systems and economic systems
    \item Understand human rights and social justice
    \item Apply historical thinking to contemporary issues
\end{itemize}

\section{Units of Study}

\textbf{Major topics include:}
\begin{itemize}
    \item World Geography and Regions
    \item Medieval Europe
    \item African and Asian Civilizations
    \item Global Trade and Exploration
    \item Enlightenment and Revolution
\end{itemize}

\section{Assessment and Grading}

\begin{center}
\begin{tabular}{|l|c|}
\hline
\textbf{Assessment Category} & \textbf{Percentage} \\
\hline
Historical Understanding & 30\% \\
\hline
Analysis and Critical Thinking & 25\% \\
\hline
Projects and Presentations & 20\% \\
\hline
Geography and Map Skills & 15\% \\
\hline
Participation & 10\% \\
\hline
\end{tabular}
\end{center}

\section{Units of Study}

\subsection{Unit 1: Geography and the Global World (5–6 weeks)}

Topics:
\begin{itemize}
\item Using maps, globes, and geographic tools
\item Physical and human geography concepts
\item Regions and cultural landscapes
\item Human–environment interaction
\item Migration, trade, and diffusion of culture
\end{itemize}

\subsection{Unit 2: Origins and Development of Civilizations (6–7 weeks)}

Topics:
\begin{itemize}
\item Transition from hunter-gatherer to agricultural societies
\item Early river valley civilizations
\item Development of writing, government, and religion
\item Trade networks and cultural exchange
\item Legacies of early civilizations
\end{itemize}

\subsection{Unit 3: Classical Civilizations and Empires (6–7 weeks)}

Topics:
\begin{itemize}
\item Classical Greece and Rome (politics, culture, philosophy)
\item Empires in Asia, Africa, and the Americas
\item Expansion, conquest, and governance
\item Cultural achievements and technological advances
\item Long-term influences on later societies
\end{itemize}

\subsection{Unit 4: World Religions and Belief Systems (5–6 weeks)}

Topics:
\begin{itemize}
\item Origins and core beliefs of major world religions
\item Practices, holidays, and traditions
\item Role of religion in societies
\item Religious art, architecture, and texts
\item Religious tolerance and conflict
\end{itemize}

\subsection{Unit 5: The Middle Ages and Global Interactions (6–7 weeks)}

Topics:
\begin{itemize}
\item Feudalism and manorialism in Europe
\item Byzantine Empire and Eastern Christianity
\item Islamic Golden Age and its contributions
\item African kingdoms and trade routes
\item Asia’s empires and cultural developments
\item Crusades and cross-cultural contact
\end{itemize}

\subsection{Unit 6: Renaissance, Reformation, and Early Modern World (6–7 weeks)}

Topics:
\begin{itemize}
\item Renaissance humanism, art, and learning
\item Reformation and religious change
\item Scientific Revolution foundations
\item Exploration and global exchange
\item Economic changes and early capitalism
\end{itemize}


\chapter{Grade 8 English/Language Arts}

\section{Course Overview}

\textbf{Grade Level:} Grade 8\\
\textbf{Duration:} Full Year (36 weeks)\\
\textbf{Prerequisites:} Grade 7 or demonstrated readiness

\section{Course Description}

Grade 8 English/Language Arts prepares students for high school English. Students analyze sophisticated texts, develop advanced writing skills, understand author's craft, and engage in academic discourse.

\section{Course Goals}

Upon completion of this course, students will be able to:

\begin{itemize}
    \item Analyze literature for complex themes and character development
    \item Determine central ideas and analyze supporting evidence
    \item Evaluate arguments and identify logical fallacies
    \item Write analytical essays with sophisticated evidence
    \item Write persuasive essays with compelling arguments
    \item Demonstrate command of standard English conventions
    \item Engage in substantive academic discussions
    \item Use precise and academic language
\end{itemize}

\section{Assessment and Grading}

\begin{center}
\begin{tabular}{|l|c|}
\hline
\textbf{Assessment Category} & \textbf{Percentage} \\
\hline
Literary Analysis and Reading & 35\% \\
\hline
Writing and Composition & 40\% \\
\hline
Grammar and Conventions & 15\% \\
\hline
Participation and Discussion & 10\% \\
\hline
\end{tabular}
\end{center}

\section{Units of Study}

\subsection{Unit 1: Analyzing Literature and Complex Texts (10-12 weeks)}

Topics:
\begin{itemize}
\item Analyzing character complexity and development
\item Understanding archetypal characters and themes
\item Analyzing plot structures and narrative techniques
\item Understanding multiple narrators and perspectives
\item Analyzing figurative language, symbolism, and literary devices
\item Making inferences and analyzing implied meanings
\item Understanding author's purpose, worldview, and craft
\item Comparing texts across cultures, genres, and time periods
\item Understanding genre characteristics and conventions
\item Analyzing dialogue and characterization techniques
\end{itemize}

\subsection{Unit 2: Non-fiction and Informational Reading (8-10 weeks)}

Topics:
\begin{itemize}
\item Analyzing complex main ideas and supporting evidence
\item Understanding text structure and organizational patterns
\item Analyzing author's perspective, bias, and arguments
\item Evaluating validity and credibility of sources
\item Analyzing rhetorical devices and persuasive techniques
\item Understanding purpose and audience
\item Synthesizing information from multiple complex sources
\item Distinguishing fact, opinion, and inference
\item Analyzing media literacy and propaganda techniques
\item Evaluating counterarguments and limitations
\end{itemize}

\subsection{Unit 3: Argumentative Writing and Debate (Ongoing)}

Duration: Integrated throughout the year

Topics:
\begin{itemize}
\item Developing nuanced claims and sophisticated thesis statements
\item Constructing logical, valid arguments
\item Supporting claims with substantial, relevant evidence
\item Acknowledging and refuting counterarguments
\item Using deductive and inductive reasoning
\item Understanding logical fallacies
\item Organizing complex arguments
\item Using credible sources and proper citations
\item Revising for logic, clarity, and persuasiveness
\item Participating in debates and collaborative argumentation
\end{itemize}

\subsection{Unit 4: Research and Informative Writing (Ongoing)}

Duration: Integrated throughout the year

Topics:
\begin{itemize}
\item Conducting extensive research using multiple sources
\item Evaluating source credibility and bias
\item Taking effective notes and organizing research
\item Developing clear central ideas
\item Integrating evidence with analysis
\item Properly citing sources in various formats
\item Avoiding plagiarism
\item Writing research papers and reports
\item Creating annotated bibliographies
\item Revising for depth and accuracy
\end{itemize}

\subsection{Unit 5: Creative and Reflective Writing (Ongoing)}

Duration: Integrated throughout the year

Topics:
\begin{itemize}
\item Writing complex narratives with multiple plot elements
\item Developing authentic characters with internal conflict
\item Using sophisticated dialogue and description
\item Employing literary devices in writing
\item Writing reflective essays
\item Understanding voice, style, and audience awareness
\item Experimenting with genre and form
\item Using revision strategies
\item Publishing for authentic audiences
\end{itemize}

\subsection{Unit 6: Presentation and Communication (Ongoing)}

Duration: Integrated throughout the year

Topics:
\begin{itemize}
\item Delivering formal presentations and speeches
\item Creating multimedia presentations
\item Engaging in academic discourse
\item Evaluating arguments and evidence in presentations
\item Using rhetorical strategies effectively
\item Practicing active listening and note-taking
\item Providing and receiving constructive critique
\item Using standard English in formal settings
\item Adapting communication for different audiences
\item Understanding non-verbal communication
\end{itemize}

\subsection{Unit 7: Vocabulary and Advanced Language Use (Ongoing)}

Duration: Integrated throughout the year

Topics:
\begin{itemize}
\item Using context clues and academic vocabulary
\item Understanding etymology and word origins
\item Analyzing word relationships and connotations
\item Building discipline-specific vocabulary
\item Understanding nuances in word meanings
\item Using reference materials strategically
\item Recognizing jargon and technical language
\item Understanding figurative language and expressions
\item Building vocabulary across disciplines
\item Using words precisely in writing and speaking
\end{itemize}

\chapter{Grade 8 Mathematics}

\section{Course Overview}

\textbf{Grade Level:} Grade 8\\
\textbf{Duration:} Full Year (36 weeks)\\
\textbf{Prerequisites:} Grade 7 or demonstrated readiness

\section{Course Description}

Grade 8 Mathematics introduces linear relationships, develops understanding of irrational numbers, explores systems of equations, and applies algebraic thinking to geometric properties and statistical relationships.

\section{Course Goals}

Upon completion of this course, students will be able to:

\begin{itemize}
    \item Understand and graph linear relationships
    \item Solve systems of linear equations
    \item Understand exponents and scientific notation
    \item Understand and identify irrational numbers
    \item Apply the Pythagorean theorem
    \item Analyze and describe transformations
    \item Understand volume and surface area
    \item Interpret statistical relationships and trends
\end{itemize}

\section{Units of Study}

\textbf{Major topics include:}
\begin{itemize}
    \item Rational and Irrational Numbers
    \item Exponents and Scientific Notation
    \item Functions and Linear Equations
    \item Systems of Equations
    \item Geometry and Transformations
    \item Statistics and Functions
\end{itemize}

\section{Assessment and Grading}

\begin{center}
\begin{tabular}{|l|c|}
\hline
\textbf{Assessment Category} & \textbf{Percentage} \\
\hline
Computational Skills & 30\% \\
\hline
Algebraic Reasoning | 35\% \\
\hline
Geometry and Statistics & 20\% \\
\hline
Participation & 10\% \\
\hline
Daily Work & 5\% \\
\hline
\end{tabular}
\end{center}

\section{Units of Study}

\subsection{Unit 1: Real Numbers and Exponents (4–5 weeks)}

Topics:
\begin{itemize}
\item Rational and irrational numbers
\item Approximating irrational numbers
\item Properties of integer exponents
\item Scientific notation and operations in scientific notation
\item Real-world applications of powers of ten
\end{itemize}

\subsection{Unit 2: Expressions and Equations (6–7 weeks)}

Topics:
\begin{itemize}
\item Writing and simplifying algebraic expressions
\item Solving multi-step linear equations
\item Solving equations with variables on both sides
\item Using properties of equality
\item Solving real-world problems with linear equations
\end{itemize}

\subsection{Unit 3: Linear Relationships and Functions (6–7 weeks)}

Topics:
\begin{itemize}
\item Understanding rate of change and slope
\item Interpreting and graphing linear functions
\item Writing equations in slope–intercept form
\item Graphing from equations, tables, and situations
\item Comparing proportional and non-proportional relationships
\item Interpreting intercepts in context
\end{itemize}

\subsection{Unit 4: Systems of Linear Equations (5–6 weeks)}

Topics:
\begin{itemize}
\item Understanding solutions of systems
\item Solving systems by graphing
\item Solving systems by substitution and elimination
\item Interpreting solutions in context
\item Modeling real-world problems with systems
\end{itemize}

\subsection{Unit 5: Geometry – Congruence, Similarity, and Transformations (6–7 weeks)}

Topics:
\begin{itemize}
\item Translations, reflections, rotations, and dilations
\item Understanding congruence through rigid motions
\item Understanding similarity through dilations
\item Using transformations to justify congruence and similarity
\item Solving problems with similar figures
\end{itemize}

\subsection{Unit 6: Geometry – Pythagorean Theorem and Distance (5–6 weeks)}

Topics:
\begin{itemize}
\item Understanding and applying the Pythagorean Theorem
\item Finding unknown side lengths in right triangles
\item Distance on the coordinate plane
\item Real-world applications of the Pythagorean Theorem
\item Exploring square roots in geometric contexts
\end{itemize}

\subsection{Unit 7: Volume of Solids (4–5 weeks)}

Topics:
\begin{itemize}
\item Volume of cylinders, cones, and spheres
\item Comparing volumes of different solids
\item Solving real-world volume problems
\item Using formulas and approximations (pi)
\end{itemize}

\subsection{Unit 8: Statistics – Bivariate Data (4–5 weeks)}

Topics:
\begin{itemize}
\item Constructing and interpreting scatter plots
\item Describing patterns of association (positive, negative, none)
\item Fitting lines informally
\item Using lines of best fit to make predictions
\item Interpreting slope and intercept in context
\end{itemize}


\chapter{Grade 8 Science}

\section{Course Overview}

\textbf{Grade Level:} Grade 8\\
\textbf{Duration:} Full Year (36 weeks)\\
\textbf{Prerequisites:} Grade 7 or demonstrated readiness

\section{Course Description}

Grade 8 Science develops deeper understanding of scientific concepts through inquiry. Students design and conduct investigations, analyze and interpret data, and develop explanations for natural phenomena.

\section{Course Goals}

Upon completion of this course, students will be able to:

\begin{itemize}
    \item Design, conduct, and refine investigations
    \item Analyze and interpret complex data
    \item Develop causal explanations
    \item Understand structure and properties of matter
    \item Understand energy transfer and transformation
    \item Understand forces and motion in depth
    \item Understand heredity and evolution
    \item Understand earth systems and human impacts
\end{itemize}

\section{Units of Study}

\textbf{Major topics include:}
\begin{itemize}
    \item Structure and Properties of Matter
    \item Chemical Reactions
    \item Energy and Energy Transfer
    \item Forces and Motion
    \item Heredity and Evolution
    \item Earth Systems and Resources
\end{itemize}

\section{Assessment and Grading}

\begin{center}
\begin{tabular}{|l|c|}
\hline
\textbf{Assessment Category} & \textbf{Percentage} \\
\hline
Investigations and Labs & 40\% \\
\hline
Data Analysis and Conclusions & 25\% \\
\hline
Scientific Reasoning | 20\% \\
\hline
Participation and Safety & 10\% \\
\hline
Science Journals & 5\% \\
\hline
\end{tabular}
\end{center}

\section{Units of Study}

\subsection{Unit 1: Matter, Atoms, and the Periodic Table (6–7 weeks)}

Topics:
\begin{itemize}
\item Atomic theory and structure (protons, neutrons, electrons)
\item Elements, compounds, and mixtures
\item Using the periodic table (groups, periods, properties)
\item Physical vs. chemical properties and changes
\item Conservation of mass in reactions
\end{itemize}

\subsection{Unit 2: Chemical Reactions and Energy (6–7 weeks)}

Topics:
\begin{itemize}
\item Evidence of chemical reactions
\item Reactants, products, and balancing simple equations (conceptual)
\item Endothermic and exothermic processes
\item Law of conservation of mass and energy in reactions
\item Real-world applications of chemical reactions
\end{itemize}

\subsection{Unit 3: Forces, Motion, and Newton’s Laws (6–7 weeks)}

Topics:
\begin{itemize}
\item Describing motion with speed, velocity, and acceleration
\item Newton’s three laws of motion
\item Net force, mass, and acceleration relationships
\item Friction, air resistance, and normal force
\item Interpreting motion graphs
\end{itemize}

\subsection{Unit 4: Energy, Work, and Simple Machines (5–6 weeks)}

Topics:
\begin{itemize}
\item Potential and kinetic energy
\item Energy transformations in systems
\item Work, power, and mechanical advantage
\item Simple machines and compound machines
\item Efficiency and energy loss (qualitative)
\end{itemize}

\subsection{Unit 5: Waves, Electromagnetic Spectrum, and Technology (6–7 weeks)}

Topics:
\begin{itemize}
\item Mechanical vs. electromagnetic waves
\item Wave properties (frequency, wavelength, amplitude, speed)
\item Sound waves and applications
\item Light behavior (reflection, refraction, absorption)
\item Electromagnetic spectrum and uses
\item Wave-based technologies (communication, medical imaging)
\end{itemize}

\subsection{Unit 6: Earth’s Systems and Climate (6–7 weeks)}

Topics:
\begin{itemize}
\item Geosphere, hydrosphere, atmosphere, and biosphere
\item Rock cycle and plate tectonics
\item Weather vs. climate
\item Factors that influence climate
\item Human impact on Earth systems
\item Natural resources and sustainability
\end{itemize}

\subsection{Unit 7: Space Systems (5–6 weeks)}

Topics:
\begin{itemize}
\item Structure of the solar system
\item Gravity and orbital motion
\item Earth’s rotation, revolution, seasons, and tides
\item Phases of the moon and eclipses
\item Stars, galaxies, and the scale of the universe
\item Space exploration and technology
\end{itemize}


\chapter{Grade 8 Social Studies}

\section{Course Overview}

\textbf{Grade Level:} Grade 8\\
\textbf{Duration:} Full Year (36 weeks)\\
\textbf{Prerequisites:} Grade 7 or demonstrated readiness

\section{Course Description}

Grade 8 Social Studies typically focuses on American history from colonization through Reconstruction. Students analyze historical events, understand causes and effects, and develop historical thinking skills.

\section{Course Goals}

Upon completion of this course, students will be able to:

\begin{itemize}
    \item Understand major periods of American history
    \item Analyze causes and effects of historical events
    \item Understand the development of American government
    \item Understand historical perspectives and interpretations
    \item Analyze the impact of historical events on society
    \item Understand the expansion and diversity of America
    \item Apply historical thinking to contemporary issues
    \item Understand the continuing relevance of historical themes
\end{itemize}

\section{Units of Study}

\textbf{Major topics include:}
\begin{itemize}
    \item Native American Civilizations
    \item European Exploration and Colonization
    \item Colonial America
    \item American Revolution
    \item Formation of Government
    \item Early Republic
    \item Westward Expansion
    \item Civil War and Reconstruction
\end{itemize}

\section{Assessment and Grading}

\begin{center}
\begin{tabular}{|l|c|}
\hline
\textbf{Assessment Category} & \textbf{Percentage} \\
\hline
Historical Understanding & 30\% \\
\hline
Analysis and Critical Thinking & 25\% \\
\hline
Projects and Presentations & 20\% \\
\hline
Maps and Geography Skills & 15\% \\
\hline
Participation & 10\% \\
\hline
\end{tabular}
\end{center}

\section{Units of Study}

\subsection{Unit 1: Foundations of the United States (6–7 weeks)}

Topics:
\begin{itemize}
\item Colonial America and regional differences
\item Causes of the American Revolution
\item Declaration of Independence and key ideas
\item Revolutionary War and outcomes
\item Articles of Confederation and need for a new constitution
\end{itemize}

\subsection{Unit 2: Constitution and Government (6–7 weeks)}

Topics:
\begin{itemize}
\item Constitutional Convention and compromises
\item Structure and principles of the Constitution
\item Federalism and separation of powers
\item Checks and balances
\item Bill of Rights and individual liberties
\item Role of citizens in a democracy
\end{itemize}

\subsection{Unit 3: Expansion and Reform (6–7 weeks)}

Topics:
\begin{itemize}
\item Early presidential administrations
\item Westward expansion and Manifest Destiny
\item Impact on Native Americans
\item Jacksonian democracy and political changes
\item Reform movements (abolition, women’s rights, temperance)
\end{itemize}

\subsection{Unit 4: Sectionalism, Civil War, and Reconstruction (6–7 weeks)}

Topics:
\begin{itemize}
\item Causes of sectional tension (slavery, economy, states’ rights)
\item Key events leading to the Civil War
\item Major battles and turning points
\item Emancipation Proclamation and outcomes
\item Reconstruction plans and amendments
\item Lasting impacts of Reconstruction
\end{itemize}

\subsection{Unit 5: Industrialization, Immigration, and Urbanization (6–7 weeks)}

Topics:
\begin{itemize}
\item Growth of industry and big business
\item Technological innovation and economic change
\item Immigration waves and cultural diversity
\item Urban growth and living conditions
\item Labor movements and reforms
\end{itemize}

\subsection{Unit 6: The United States as a World Power (5–6 weeks)}

Topics:
\begin{itemize}
\item Imperialism and expansion overseas
\item Spanish–American War and consequences
\item World War I causes and U.S. involvement
\item Treaty of Versailles and League of Nations debate
\item Changing role of the U.S. in world affairs
\end{itemize}

\subsection{Unit 7: Modern United States (5–6 weeks)}

Topics:
\begin{itemize}
\item Major events of the 20th and 21st centuries (overview)
\item Civil Rights Movement and expansion of rights
\item Social and technological change
\item Contemporary issues in government and society
\item Role of the United States in a globalized world
\end{itemize}


\chapter{Grade 7-8 Overview}

This chapter provides an overview of middle school (Grades 7-8) curriculum, including:

\begin{itemize}
    \item Developmental characteristics of middle school students
    \item Transition from elementary to middle school
    \item Academic expectations and rigor
    \item Social-emotional learning integration
    \item Interdisciplinary connections
    \item Technology integration
    \item Differentiation and support strategies
\end{itemize}

\clearpage

\appendix

\chapter{Standards Alignment}

This appendix provides detailed alignment with:

\begin{itemize}
    \item Common Core State Standards (CCSS)
    \item Next Generation Science Standards (NGSS)
    \item National Council for the Social Studies (NCSS) Standards
    \item National Standards for Arts Education
    \item SHAPE America Physical Education Standards
    \item National Health Education Standards
\end{itemize}

\chapter{Differentiation Strategies}

This appendix outlines strategies for:

\begin{itemize}
    \item Differentiation for varied learners
    \item Support for English Language Learners (ELL)
    \item Accommodations for students with special needs
    \item Enrichment for advanced learners
    \item Universal Design for Learning (UDL) principles
    \item Response to Intervention (RTI) frameworks
\end{itemize}

\chapter{Assessment Resources}

This appendix includes:

\begin{itemize}
    \item Assessment rubrics and tools
    \item Formative assessment strategies
    \item Data collection methods
    \item Standards-based grading guidelines
    \item Performance task examples
    \item Student progress monitoring tools
\end{itemize}

\clearpage

\chapter*{Conclusion}
\addcontentsline{toc}{chapter}{Conclusion}

This comprehensive curriculum guide provides institutional-grade documentation for K-8 education. The framework ensures coherence across grade levels while allowing flexibility for local adaptation and differentiation.

The curricula are designed to:

\begin{itemize}
    \item Meet national and state standards
    \item Support students with diverse learning needs
    \item Promote critical thinking and problem-solving
    \item Develop engaged and informed citizens
    \item Prepare students for future success
\end{itemize}

Successful implementation requires ongoing professional development, collaboration among teachers, and attention to the individual needs of all learners.

\end{document}
